\chapter{文本分类、匹配与组织}
\label{chap:nlp-classification-matching-organization}

从通知中抽出活动、日期和报名条件以后，还有另一组问题没有回答：这份通知应分到哪些栏目？两份通知是不是在说同一件事？数千份尚无目录的材料可以怎样分组？第~\ref{chap:nlp-information-extraction}章关注原文中的对象、关系与记录，本章转而考察三类对象：单篇文本的属性、两段文本之间的关系，以及文本集合的组织结构。处理对象变了，监督从哪里来、结果怎样验收也随之改变。

沿用本部分的教学通知：活动编号为\code{ML-0418}的“机器学习读书会”于2025年4月15日发布更正：“机器学习读书会原定于4月17日举行，现改至4月18日（周五）14:00，地点为海棠楼201。校内师生可直接参加，校外来宾须在4月16日18:00前提交报名信息。”这份材料既有时间安排，又有报名说明；删去报名条件的短句仍与它主题相关，却未必是等义复述；把它分进“报名”文档簇，也不意味着每句话都讨论报名。本章的通知、评论、向量、概率和计数都是教学构造，不是模型实测结果。

表~\ref{tab:nlp-classification-task-meaning}把三类任务放在同一材料上比较。属性判定有逐篇标注，按标签评价；关系判定的监督是成对或有序的标注，评价要看方向与候选集合；集合组织通常没有标注，结果的含义由所用模型决定，验收依赖人工解释与稳定性检查。

\begin{table*}[htbp]
  \centering
  \small
  \begin{tabularx}{\linewidth}{@{}>{\raggedright\arraybackslash}p{16mm}
    >{\raggedright\arraybackslash}p{22mm}>{\raggedright\arraybackslash}X
    >{\raggedright\arraybackslash}p{23mm}>{\raggedright\arraybackslash}X
    >{\raggedright\arraybackslash}p{22mm}@{}}
    \toprule
    任务 & 输入对象 & 输出含义 & 监督依据 & 通知中的反例 & 评价口径 \\
    \midrule
    多标签分类 & 一份通知 & 同时包含哪些内容 & 栏目标签规范 & 有报名词不等于只谈报名 & 逐标签及宏、微F1 \\
    情感 & 评论及评价对象 & 对对象的极性或强度 & 指定粒度的意见 & 抱怨改期不等于反对活动 & 按对象、极性评价 \\
    立场 & 文本与指定主张 & 是否支持该主张 & 目标条件标注 & 语气消极仍可支持改期 & 按目标和标签评价 \\
    相似度 & 两段文字 & 意义接近的程度 & 连续参考评分 & 日期不同仍可能得高分 & Pearson、Spearman \\
    复述 & 两段文字 & 约定范围内是否等义 & 等义句对标注 & 删掉报名截止条件 & F1、条件反例 \\
    推断 & 前提与假设 & 前提支持、排斥或未定 & 有方向的三分类 & 有时间不推出有主讲人 & 三分类与方向测试 \\
    相关性 & 查询与候选 & 对当前需求是否有用 & 查询条件相关性 & 讲读书会却不含新时间 & Recall@k、排序指标 \\
    对齐 & 两组文本单元 & 哪些单元相互对应 & 平行材料或对应标注 & 译文把一句拆成两句 & 覆盖率与配对正确性 \\
    主题 & 文档集合 & 共现结构及混合权重 & 无标签词语观测 & 主题比例不是栏目概率 & 连贯性及路线专属指标 \\
    聚类 & 文档集合 & 在给定距离下的组别 & 表示和聚类准则 & 同簇通知可有不同日期 & 稳定性、轮廓及业务核验 \\
    \bottomrule
  \end{tabularx}
  \caption{属性、关系与组织结果的区别。各行围绕同一活动材料改变输入和判断问题；共享分类头或向量表示，不会使这些标签相互等价。F1的统计单位、相关性的候选集合和组织结果的解释条件都需随任务记录。}
  \label{tab:nlp-classification-task-meaning}
\end{table*}

三类任务各有自己的困难。属性判定要面对标签语义与粒度：一篇文本可带多个标签，标签之间有层次，商品类目这类任务的标签数达到数十万；长文档中决定类别的只有少数句子；情感与立场依赖评价对象或目标主张，同一句话对不同对象结论不同；领域迁移和新增标签使训练好的分类器失效；来源判定还要面对改写与对抗，生成文本与人写文本在统计上越来越接近。关系判定的核心是权衡：两段文本在编码时交互越充分，越能分辨角色交换与条件遗漏，可每个候选对都要完整计算，候选规模一大就无法承受；关系的方向（蕴含有方向，相似没有）也必须进入模型和评价。集合组织没有标注，主题数、簇数需要选择，结果的解释与评价都缺少标准答案。

方法谱系沿三条线展开。第~\ref{sec:nlp-classification-attributes}节的分类从词袋线性模型经局部卷积与层次注意力网络，走到预训练编码器与大语言模型提示，每一步减少一部分特征设计或标注需求。第~\ref{sec:nlp-classification-matching}节按两段文本何时交互区分表示型、交互型与后期交互架构，相似度、复述、推断与相关性模型都在其中定位。第~\ref{sec:nlp-classification-organization}节从矩阵分解、概率与神经主题模型走到表示聚类和关键词组织。

本章还复用第~\ref{sec:nlp-analysis-paradigms}节的结构预测范式。方面级情感分析的统一抽取方法分别对应约束生成、跨度枚举与词对表三种范式；关键短语抽取在有参考短语时转为序列标注或跨度分类。这些地方只写任务特有的输出设计，解码与训练目标回指第~\ref{chap:nlp-text-analysis}章。朴素贝叶斯、逻辑回归与支持向量机、潜在语义分析、聚类算法以及主题模型的一般理论属于前卷，各节在首次用到时说明所用结论并引用相应的章；本章只讲相关性模型本身，不讲检索系统的索引与在线架构。

\FloatBarrier
\section{文本分类与属性判定}
\label{sec:nlp-classification-attributes}

给一段文字贴标签之前，应先问标签在描述什么。“报名说明”描述内容，“通知”描述体裁，“提醒”描述用途；“不满”描述情绪，“支持改期”描述对指定方案的立场。它们可以同时成立，把所有名称放进同一个互斥类别表，会把本来可以并存的属性强迫成竞争关系。分类单位也须固定：一句表达不满的话，不一定代表整篇通知的语气，更不能代表转发者的态度。

本节先以内容类别展开分类模型的谱系，再处理多标签与极端多标签；情感、立场与文本来源三类属性各需一个额外条件：评价对象、目标主张，或可访问的来源模型。

\subsection{内容类别与用途判定}
\label{sec:nlp-classification-content}

\term{单标签分类}（single-label classification）要求一个样本只选一个类别；\term{多标签分类}（multi-label classification）允许多个属性同时为真。\term{层次分类}（hierarchical classification）再加上类别之间的包含关系，例如“活动通知”之下再分“时间安排”和“报名说明”；这种包含属于标注规范，不由词义近似自动推出。业务还应约定“其他”“未知”和“证据不足”：它们分别表示已知的剩余类别、超出类别表的输入，以及当前文本不足以判定，三者不宜合并。

设训练样本为$(x_i,\bm y_i)$，$x_i$是文档，$\bm y_i\in\{0,1\}^{K}$记录$K$个标签。文档表示$\bm h_i\in\mathbb R^d$经过$\bm z_i=\bm W\bm h_i+\bm b$得到标签得分。互斥类别使用
\begin{equation}
  \label{eq:nlp-classification-softmax}
  p(y=k\mid x)=\operatorname{softmax}(\bm z)_k
  =\frac{\exp(z_k)}{\sum_{\ell=1}^{K}\exp(z_\ell)},
\end{equation}
分母遍历全部候选类别，提高一类的得分会压低其他类的概率。多标签改用逐标签的sigmoid与二元交叉熵：
\begin{equation}
  \label{eq:nlp-classification-sigmoid}
  \begin{aligned}
    p_k&=\sigma(z_k)=\frac{1}{1+\exp(-z_k)},\\
    L(x,\bm y)&=-\sum_{k=1}^{K}
    \left[y_k\log p_k+(1-y_k)\log(1-p_k)\right].
  \end{aligned}
\end{equation}
每个$p_k$估计“标签$k$存在”的条件概率，彼此不要求加和为1。生成式与判别式分类的一般理论，包括朴素贝叶斯、逻辑回归与支持向量机的定义、学习准则和求解，可参考\BookChapterRef{chap:classification}{卷四的分类章}；本节只写文本带来的特殊问题：特征怎样从词序列中构造，以及长度、词序和标签规模怎样影响模型设计。

文本分类的挑战可以归为四项。词特征稀疏而高维，一个类别的证据分散在成千上万个词上；词序与局部短语（“无需报名”与“需要报名”）改变含义；长文档里关键信息只占少数句子，截断或平均都会稀释它；标注不足，类别表还会新增标签。表~\ref{tab:nlp-classification-content-methods}按这四项列出本小节的方法谱系。

\begin{table*}[htbp]
  \centering
  \small
  \begin{tabularx}{\linewidth}{@{}>{\raggedright\arraybackslash}p{24mm}
    >{\raggedright\arraybackslash}X>{\raggedright\arraybackslash}X
    >{\raggedright\arraybackslash}X@{}}
    \toprule
    方法 & 解决的问题 & 核心机制 & 新代价或局限 \\
    \midrule
    多项式朴素贝叶斯 & 样本少时给出可检查的基线 & 类内词分布与平滑计数，按后验选类 & 条件独立过强，重复词造成过度自信 \\
    逻辑回归、SVM & 直接学判别边界，容许相关特征 & TF-IDF与$n$-gram稀疏特征上的线性得分 & 特征工程依赖人工，同义词不共享统计 \\
    fastText & 让相似词与短语共享参数且保持线性速度 & 词与哈希$n$-gram嵌入平均后接线性头 & 平均丢失长距离结构，哈希碰撞 \\
    TextCNN & 学习不必枚举的局部短语模式 & 多宽度卷积与时间维最大池化 & 只看窗口内，跨句条件难以组合 \\
    层次注意力网络 & 长文档中关键句被平均或截断 & 词级与句级两层双向GRU加注意力汇聚 & 顺序计算慢，注意力权重不等于证据 \\
    BERT类编码器 & 否定、对象和条件的上下文交互 & 预训练双向编码与任务头微调 & 输入长度受限，每任务需标注与微调 \\
    大语言模型提示 & 无标注或新标签时仍需分类 & 标签词条件概率或受约束生成 & 每次调用成本高，对提示与示例敏感 \\
    \bottomrule
  \end{tabularx}
  \caption{内容分类方法谱系。自上而下，前一行的局限大致是后一行的动机；代价一栏说明每一步换来的新问题，不表示后面的方法在所有数据上都更好。}
  \label{tab:nlp-classification-content-methods}
\end{table*}

\subsubsection{多项式朴素贝叶斯文本分类}
\label{sec:nlp-classification-naive-bayes}

监督样本不多、又希望检查每个词怎样影响判断时，可以从类内词频建立生成式基线。文本特有的第一个选择是\emph{事件模型}：把文档看成从词表中有放回地抽出的一串词（多项式模型），还是看成每个词“出现或不出现”的一组二元事件（伯努利模型）\scite{mccallum1998event}。固定训练词表$V$，令$x_v$为文档中词$v$的出现次数，$c_{k,v}$为类别$k$训练文档中$v$的总计数，$n_k$为该类文档数。\term{多项式朴素贝叶斯}（multinomial naive Bayes）用加性平滑估计
\[
  \widehat\pi_k=\frac{n_k}{n},\qquad
  \widehat\phi_{k,v}
  =\frac{c_{k,v}+\lambda}{\sum_{u\in V}c_{k,u}+\lambda|V|},
  \quad \lambda>0,
\]
并按$s_k=\log\widehat\pi_k+\sum_vx_v\log\widehat\phi_{k,v}$选类；对所有类别相同的多项式系数已约去。条件独立假设作用于给定类别后的词项事件，训练只需保存先验和类内词分布。

考虑区分“活动通知”和“其他材料”的教学词表$\{\text{读书会},\text{报名},\text{比赛}\}$。两类训练文档数为3、2，词计数分别为$(4,3,1)$、$(1,1,6)$。取$\lambda=1$，通知类词概率为$(5,4,2)/11$，其他类为$(2,2,7)/11$。待分类通知的计数为$(1,2,0)$，
\[
  \begin{aligned}
    s_{\text{通知}}&=\log0.6+\log\tfrac{5}{11}+2\log\tfrac{4}{11}\approx-3.3225,\\
    s_{\text{其他}}&=\log0.4+3\log\tfrac{2}{11}\approx-6.0305,
  \end{aligned}
\]
通知类领先约2.708，两类归一化后的后验约为0.9375。

伯努利模型改记“含有$v$的文档数”。设通知类3篇中含三词的篇数为$(3,2,1)$，其他类2篇中为$(1,1,2)$，平滑后出现概率为$(\tfrac45,\tfrac35,\tfrac25)$与$(\tfrac24,\tfrac24,\tfrac34)$。同一通知的出现向量为$(1,1,0)$，似然还要乘上“比赛未出现”的概率：通知类为$0.6\times0.8\times0.6\times0.6=0.1728$，其他类为$0.4\times0.5\times0.5\times0.25=0.025$，后验约0.8736。两种模型的差别正落在文本特征上：多项式模型让重复两次的“报名”贡献两次证据，伯努利模型只记出现一次，却把缺席词也当作证据。短文本、词表较小时伯努利模型往往够用；长文档里词频携带信息，多项式模型更常用。

平滑使已在词表中、但训练时未在某类出现的词仍有正概率；词表之外的词须预先约定忽略或映射到未知词项。重复模板会让多项式模型累积过强的置信；词袋也无法表达否定范围，“无需报名”和“需要报名”只差一个计数。

\subsubsection{线性文本分类：逻辑回归与SVM}
\label{sec:nlp-classification-linear}

朴素贝叶斯要求词在类内独立，而“报名”与“截止”这类常同现的词会被重复计权。判别式线性模型不再为词建模，直接学习哪些特征区分类别，相关特征之间的权重由训练自动分配。文本上的工作主要在特征：词与相邻词对的$n$-gram计数，经词频—逆文档频率（term frequency–inverse document frequency, TF-IDF）加权。本章采用$\operatorname{idf}_v=\log((n+1)/(\operatorname{df}_v+1))+1$，$\operatorname{df}_v$为训练语料中含$v$的文档数，词频乘该权重后作$L_2$归一化。词表、低频过滤和IDF只在训练语料上拟合，验证文档不参与这些统计。

保留“读书会、报名、比赛”三维，通知向量为$\bm x=(0.6,0.8,0)^\top$。逻辑回归头取$\bm w=(1,1,-1)^\top$、$b=-0.5$，得分0.9，通知概率$\sigma(0.9)\approx0.7109$，参考标签为1时损失$-\log0.7109\approx0.3412$。线性SVM取$\bm w=(0.5,1,-1)^\top$、$b=-0.2$，得分也为0.9，对$y=+1$的铰链损失$\max(0,1-ys)=0.1$。两个0.9都决定了边界的同一侧，但SVM得分是到边界的带符号间隔，需要另拟合校准映射才能当概率用。

稀疏表示使得分只需累加非零特征，训练与预测都与文档中不同词的个数成正比。$L_1$正则化可以令多数$n$-gram权重为零，起特征选择作用；正类稀少时，可对损失加类别权重，阈值仍在符合目标分布的验证集上确定。这类模型在主题分类上至今是很强的基线，它的局限也来自特征：“报名”与“登记”是两个互不相关的维度，未在训练集出现的短语没有权重。

\subsubsection{fastText：词与n-gram表示聚合}
\label{sec:nlp-classification-fasttext}

让同义词共享统计的直接办法，是给每个词和短语一个可学习的低维向量。fastText分类器把文档的特征多重集$F(x)$（词与词$n$-gram）中的向量取平均，再接线性分类：
\[
  \bm h(x)=\frac{1}{|F(x)|}\sum_{f\in F(x)}\bm e_f\in\mathbb R^{d},\qquad
  \bm z=\bm W\bm h+\bm b.
\]
交叉熵梯度同时更新特征嵌入表与分类矩阵。原方法有两处为速度而设：$n$-gram数量随语料膨胀，便用哈希函数把它们映射到固定数目的桶中共享向量；类别很多时，输出层改用第~\ref{sec:nlp-classical-hierarchical-softmax}小节的层次softmax，把每个样本的输出计算从$O(Kd)$降到$O(d\log K)$\scite{nlp-joulin2017fasttext}。

fastText可以看作低秩的逻辑回归：嵌入维度$d$较小时，它把逻辑回归的稀疏权重分解为“特征嵌入$\times$分类矩阵”，相似的词通过相近的嵌入共享统计；去掉$n$-gram，它就是词袋的平均嵌入分类器。

$n$-gram特征的作用可以用两句话检验：“先／报名／再／参加”与“先／参加／再／报名”。只用词特征时两份特征多重集相同，任何参数下得分都相同。加入相邻二元短语后，第一句多出“先报名、报名再、再参加”，第二句多出“先参加、参加再、再报名”，两份多重集大小都为7。设共同四个词向量之和为$(2,1)$，第一句三个短语向量之和为$(2,0)$，第二句为$(0,2)$，均值表示分别为$(4/7,1/7)$与$(2/7,3/7)$；若“报名先行”标签头为$\bm w=(2,-1)$、$b=0$，两句得分为1与$1/7$，sigmoid约为0.7311与0.5357，顺序信息由短语特征带入。这里的向量是为说明信息路径而给定的教学数值。

$n$-gram只保留窗口内的顺序，平均仍会抹掉长距离结构，哈希碰撞也可能把无关短语混在同一向量里；换来的是速度，训练和预测只需查表、求和与一次矩阵乘法。

\subsubsection{TextCNN：局部模式与池化}
\label{sec:nlp-classification-textcnn}

$n$-gram表需要枚举短语，近义表达“提交报名”“递交登记”仍是不同的条目。TextCNN让滤波器在连续词向量窗口上学习局部模式，相近的短语因词向量相近而得到相近的响应\scite{nlp-kim2014textcnn}。设输入矩阵$\bm X\in\mathbb R^{n\times d}$的第$j$行是词向量$\bm e_j$，宽为$r$的一个滤波器$\bm w\in\mathbb R^{rd}$在每个位置计算
\[
  a_j=f\!\left(\bm w^\top
  [\bm e_j;\ldots;\bm e_{j+r-1}]+b\right),\quad j=1,\ldots,n-r+1,\qquad
  h=\max_j a_j.
\]
分号表示拼接，$f$为ReLU。时间维最大池化把长度不一的响应序列变成一个数，并让梯度只流回被选中的窗口。原文取宽度3、4、5各100个滤波器，池化后拼成300维向量，经dropout接softmax；张量形状依次为$n\times d\to(n-r+1)\times100$（每种宽度）$\to300\to K$。它还比较了两种词向量通道：静态通道固定预训练词向量，非静态通道随任务微调，双通道则同时保留二者，使微调不至于把预训练的相似结构完全改掉。卷积与池化的通用原理可参考\BookChapterRef{chap:convolutional-networks}{卷五的卷积神经网络章}。

把“来宾须提交报名信息”切成五个教学词元，暂用一维词向量$(0,0.5,1,1,0)$。宽2滤波器取$\bm w=(1,1)$、$b=-0.5$，四个响应为$(0,1,1.5,0.5)$，最大响应1.5来自“提交／报名”；若报名头为$z=2h-1$，则$z=2$，概率0.8808。原文改为“无需提交报名信息”时，触发窗口仍在，否定词若落在窗口外，最大池化保留的仍是错误线索。

TextCNN相对fastText增加了可学习的局部组合，视野仍是固定宽度的窗口：“校内无需报名；校外须提交信息”需要按对象组织两条规则，跨句条件也超出窗口；长文档整篇卷积后取全局最大，每个滤波器只留下最强的一处响应。

\subsubsection{层次注意力网络：长文档分类}
\label{sec:nlp-classification-han}

长文档带来两个问题。其一，类别证据往往集中在少数句子里，一篇几百词的评论可能只有两句决定评分，平均嵌入把它们稀释在其余句子中，截断又可能把它们丢掉。其二，文档有“词组成句、句组成篇”的层次，平铺的词序列模型看不到句子边界。\term{层次注意力网络}（hierarchical attention network, HAN）按这两层结构分别编码与汇聚：先在每个句子内部用注意力挑出重要的词得到句向量，再在句子之间用注意力挑出重要的句子得到文档向量\scite{nlp-yang2016han}。整体结构见图~\ref{fig:nlp-classification-han}。

\begin{figure*}[tbp]
  \centering
  % 层次注意力网络：词级双向GRU与词注意力得到句向量，句级双向GRU与句注意力得到文档向量。
\begin{tikzpicture}[
  x=1mm,y=1mm,figure picture,
  cell/.style={figure node,minimum width=11mm,minimum height=5.5mm,
    inner sep=1pt,font=\figurefont\scriptsize},
  word/.style={font=\figurefont\scriptsize},
  row label/.style={font=\figurefont\scriptsize,text=FigureMuted,anchor=east},
  flow/.style={figure flow},
  att/.style={draw=FigureOutput,line width=.7pt,-{Stealth[length=1.3mm]}},
  rnn/.style={draw=FigureStroke,line width=.7pt,{Stealth[length=1.2mm]}-{Stealth[length=1.2mm]}}
]
  % 词级
  \foreach \i/\xo in {1/0,2/48}{
    \foreach \t in {1,2,3}{
      \pgfmathsetmacro{\xx}{\xo+(\t-1)*14}
      \node[word] (w\i\t) at (\xx,0) {$x_{\i,\t}$};
      \node[cell,draw=FigureInput] (h\i\t) at (\xx,10) {$\bm h_{\i,\t}$};
      \draw[flow] (w\i\t) -- (h\i\t);
    }
    \draw[rnn] (h\i1) -- (h\i2);
    \draw[rnn] (h\i2) -- (h\i3);
    \node[cell,draw=FigureModel] (s\i) at (\xo+14,24) {$\bm s_{\i}$};
    \foreach \t in {1,2,3}{
      \draw[att] (h\i\t) -- node[pos=.45,font=\figurefont\tiny,fill=white,inner sep=.5pt]
        {$\alpha_{\i,\t}$} (s\i);
    }
  }
  \node[cell,draw=FigureMuted,densely dashed] (uw) at (38,24) {查询$\bm u_w$};
  \draw[figure link,densely dashed] (uw) -- (s1);
  \draw[figure link,densely dashed] (uw) -- (s2);
  % 句级
  \node[cell,draw=FigureInput] (g1) at (14,38) {$\bm g_{1}$};
  \node[cell,draw=FigureInput] (g2) at (62,38) {$\bm g_{2}$};
  \draw[flow] (s1) -- (g1);
  \draw[flow] (s2) -- (g2);
  \draw[rnn] (g1) -- (g2);
  \node[cell,draw=FigureModel] (v) at (38,52) {$\bm v$};
  \draw[att] (g1) -- node[pos=.5,font=\figurefont\tiny,fill=white,inner sep=.5pt] {$\beta_1$} (v);
  \draw[att] (g2) -- node[pos=.5,font=\figurefont\tiny,fill=white,inner sep=.5pt] {$\beta_2$} (v);
  \node[cell,draw=FigureMuted,densely dashed] (us) at (62,52) {查询$\bm u_s$};
  \draw[figure link,densely dashed] (us) -- (v);
  \node[cell,draw=FigureOutput,minimum width=22mm] (y) at (38,64) {$\operatorname{softmax}(\bm W_c\bm v+\bm b_c)$};
  \draw[flow] (v) -- (y);
  \node[row label] at (-8,0) {词元};
  \node[row label] at (-8,10) {词级双向GRU};
  \node[row label] at (-8,24) {词注意力汇聚为句向量};
  \node[row label] at (-8,38) {句级双向GRU};
  \node[row label] at (-8,52) {句注意力汇聚为文档};
  \node[row label] at (-8,64) {分类};
\end{tikzpicture}

  \caption{层次注意力网络。下半部分在每个句子内部运算：词级双向GRU给出$\bm h_{i,t}$，以可学习查询$\bm u_w$计算词注意力$\alpha_{i,t}$，加权得句向量$\bm s_i$；上半部分在句子序列上重复同一模式，以查询$\bm u_s$计算句注意力$\beta_i$，汇聚为文档向量$\bm v$后分类。}
  \label{fig:nlp-classification-han}
\end{figure*}

设文档有$L$个句子，第$i$句有$T_i$个词，词向量为$\bm x_{i,t}\in\mathbb R^{d_e}$。词级编码用\term{门控循环单元}（gated recurrent unit, GRU）分别从左到右、从右到左读句子，拼接两方向的隐状态\cite{cho2014}；GRU用更新门决定保留多少旧状态，结构与LSTM相近而门更少。词注意力先把每个词表示投影成“重要性键”，再与一个全局可学习的查询向量比较：
\begin{equation}
  \label{eq:nlp-classification-han-word}
  \begin{aligned}
    \bm h_{i,t}&=\bigl[\overrightarrow{\mathrm{GRU}}(\bm x_{i,t});\ \overleftarrow{\mathrm{GRU}}(\bm x_{i,t})\bigr]\in\mathbb R^{2d_h},\\
    \bm u_{i,t}&=\tanh(\bm W_w\bm h_{i,t}+\bm b_w),\\
    \alpha_{i,t}&=\frac{\exp(\bm u_{i,t}^\top\bm u_w)}{\sum_{t'=1}^{T_i}\exp(\bm u_{i,t'}^\top\bm u_w)},\qquad
    \bm s_i=\sum_{t=1}^{T_i}\alpha_{i,t}\bm h_{i,t}.
  \end{aligned}
\end{equation}
句级编码在句向量序列$\bm s_1,\ldots,\bm s_L$上重复同一模式：
\begin{equation}
  \label{eq:nlp-classification-han-sentence}
  \begin{aligned}
    \bm g_i&=\bigl[\overrightarrow{\mathrm{GRU}}(\bm s_i);\ \overleftarrow{\mathrm{GRU}}(\bm s_i)\bigr],\qquad
    \bm u_i=\tanh(\bm W_s\bm g_i+\bm b_s),\\
    \beta_i&=\frac{\exp(\bm u_i^\top\bm u_s)}{\sum_{i'=1}^{L}\exp(\bm u_{i'}^\top\bm u_s)},\qquad
    \bm v=\sum_{i=1}^{L}\beta_i\bm g_i,\\
    p(y\mid x)&=\operatorname{softmax}(\bm W_c\bm v+\bm b_c).
  \end{aligned}
\end{equation}
查询$\bm u_w,\bm u_s$与其余参数一起由文档级交叉熵训练，没有词或句的重要性标注。原文取词向量200维、每方向GRU 50维，于是一个批次的张量形状依次为：输入$B\times L\times T\times200$，词级编码$B\times L\times T\times100$，词注意力汇聚后$B\times L\times100$，句级编码仍为$B\times L\times100$，句注意力汇聚后$B\times100$，输出$B\times K$。实现时把同批文档补齐到相同的$L$和$T$，补齐位置在softmax前置为$-\infty$。

用一个句级的教学数值看注意力怎样改变汇聚。设三句的句级表示为$\bm g_1=(1,0)$、$\bm g_2=(0,1)$、$\bm g_3=(1,1)$，投影后的键$\bm u_i$为$(0.2,0.1)$、$(0.9,0.8)$、$(0.1,-0.1)$，查询$\bm u_s=(1,1)$。打分为$(0.3,1.7,0)$，softmax得$\bm\beta\approx(0.173,0.700,0.128)$，文档向量$\bm v\approx(0.300,0.828)$；平均汇聚则得$(0.667,0.667)$。若第二句是“校外来宾须提前报名”，句注意力让它主导文档表示，而平均把它与另外两句等量混合。

把两层注意力换成平均，HAN退化为层次平均编码；把整篇当作一句，它就是带注意力的BiGRU。它的收益来自句子边界这一结构先验，以及不增加标注就学到的句子权重。代价是循环网络按顺序计算，长文档训练较慢；注意力权重只是汇聚系数，可帮助定位关键句，却未必是判断的充分原因。超过编码器长度上限的文档仍可沿用这种“段内编码、段间注意力汇聚”的层次思路（第~\ref{sec:nlp-classification-bert}小节）。

\subsubsection{BERT类模型：上下文编码与分类头}
\label{sec:nlp-classification-bert}

前面的模型从零学习或只用静态词向量，否定、对象和条件的相互作用要靠有限的任务标注学出。BERT类模型先在大规模无标注文本上预训练双向编码器，再取\code{[CLS]}位置的表示$\bm h_{\mathrm C}\in\mathbb R^{H}$接线性头，用式~\eqref{eq:nlp-classification-softmax}或式~\eqref{eq:nlp-classification-sigmoid}微调全部参数。分类头的结构与微调流程见第~\ref{sec:nlp-classical-bert}小节，这里不再重复，只讨论文本分类落地时的三个问题：多标签阈值、层次一致与长文档。

多标签训练与阈值选择需要分开计算。只考察“时间安排、报名说明”两个标签，通知参考值为$\bm y=(1,1)$。取$\bm h=(1,2)^\top$、$\bm b=\bm0$，分类矩阵两行为$(\log4,0)$与$(\log1.5,0)$，则$\bm p=(0.8,0.6)$，单样本损失$L=-\log0.8-\log0.6=0.7340$，$\partial L/\partial\bm z=\bm p-\bm y=(-0.2,-0.4)^\top$。冻结编码器和偏置，以学习率0.1更新$\bm W^{(1)}=\bm W^{(0)}-0.1(\bm p-\bm y)\bm h^\top$，两行分别增加$(0.02,0.04)$和$(0.04,0.08)$，概率变为约$(0.8155,0.6469)$。这一步只用训练标签，没有选择阈值。

阈值要在冻结模型后另取的验证集上选择。表~\ref{tab:nlp-classification-thresholds}给出四篇验证文档，候选阈值取0.4、0.5、0.6三档，$p_k\geq t_k$即输出标签$k$。报名标签在阈值0.4下有两次真阳性、一次假阳性、零次假阴性，F1为$0.8$；阈值0.5的F1为0.5，0.6为$2/3$，故选0.4。时间标签在0.5与0.6下都完全正确，按事先约定的“并列取较小阈值”选0.5。

\begin{table}[htbp]
  \centering
  \begin{tabular}{@{}ccccc@{}}
    \toprule
    验证文档 & 时间得分 & 报名得分 & 时间参考 & 报名参考 \\
    \midrule
    A & 0.80 & 0.60 & 1 & 1 \\
    B & 0.70 & 0.55 & 1 & 0 \\
    C & 0.40 & 0.45 & 0 & 1 \\
    D & 0.30 & 0.35 & 0 & 0 \\
    \bottomrule
  \end{tabular}
  \caption{用于选择阈值的教学验证集，与训练样本分开。B只有时间安排却获得偏高的报名分数，C的报名说明得分偏低，使两个标签需要不同的阈值。}
  \label{tab:nlp-classification-thresholds}
\end{table}

逐标签阈值$(0.5,0.4)$下，宏F1为$(1+0.8)/2=0.9$；两个标签的计数相加得$\mathrm{TP}=4,\mathrm{FP}=1,\mathrm{FN}=0$，微F1为$8/9$。宏平均给每个标签相同权重，罕见标签的错误因而更显眼；微平均给每次标签判定相同权重，由常见标签主导。若只许一个全局阈值，0.4、0.5、0.6的宏F1分别为0.8、0.75、$5/6$，按宏F1选0.6。小验证集上的最优阈值不稳定，罕见标签应报告支持数，必要时退回共享阈值。

层次标签要求输出合法：细标签“报名说明”入选而父标签“活动通知”未入选时，可按规范补入祖先，也可在训练和解码中显式约束子标签概率不超过父标签。长文档超过编码器的输入上限（BERT为512个词元）时，先按重叠窗口分段编码，再对段表示做平均、最大池化或上一小节的注意力汇聚：平均可能稀释短的报名条款，最大池化可能忽略它的适用对象，注意力汇聚需要额外训练段级参数。

固定的$K$维分类头只在已知类别之间分配概率，最高分不能说明“所有类别都不适用”。未知拒绝需要独立条件，例如最大概率低于在含未知样本的开发集上确定的门槛时输出“未知”，并单独报告未知样本的召回与误拒。类别表需要扩展时，固定分类头就要重建，这是下一小节提示式分类的出发点。

\subsubsection{大语言模型的零样本与少样本分类}
\label{sec:nlp-classification-llm}

微调分类器需要每个任务的标注数据和一次训练；类别表一变，输出层就要重建。大语言模型提供另一条路：把类别写成自然语言，让模型在条件分布中比较各标签的可能性，不训练任何分类头。

最常用的打分方式是\term{标签词}（verbalizer）：为每个类别$k$指定一个词或短语$v_k$，例如“时间安排”“报名”“其他”，把输入包进模板$\operatorname{prompt}(x)$——“通知：$x$。这则通知的主要内容是：”——再比较各标签词作为续写的条件概率\cite{nlp-schick2021pet}：
\begin{equation}
  \label{eq:nlp-classification-verbalizer}
  s_k(x)=\frac{1}{|v_k|}\sum_{j=1}^{|v_k|}\log p_{\bm\theta}\bigl(v_{k,j}\mid \operatorname{prompt}(x),v_{k,<j}\bigr),\qquad
  \hat y=\arg\max_k s_k(x).
\end{equation}
标签词可能被切成多个词元，求和会偏向短标签，因此这里按词元数平均；也可以把类别编成选项字母“A/B/C”，只比较一个词元的概率。另一种做法让模型自由生成标签名，再用受约束解码把输出限定在合法类别内，或解析后映射。打分式只需一次前向，就能得到全部类别的相对分数，便于设阈值和拒绝；生成式可以附带理由，但非法输出要单独计数，不能静默映射到最近的类别。

\term{零样本提示}（zero-shot prompting）只给任务说明和当前输入；\term{少样本提示}（few-shot prompting）在说明后附$k$个带标签的示例，二者都属于第~\ref{sec:nlp-inference-icl-shots}小节的上下文学习，参数不变，只改条件序列。示例主要约束标签空间与输出格式，也带来系统偏差：模型可能偏爱示例中多数的标签、最后一个示例的标签，或预训练中常见的标签词。

第~\ref{sec:nlp-inference-icl-sensitivity}小节的上下文校准正是针对这类标签偏好的修正：保持说明与示例不变，用“N/A”等无内容输入读出标签概率$\hat{\bm p}_{\mathrm{cf}}$，再按式~\eqref{eq:nlp-inference-calibration}以$\operatorname{diag}(\hat{\bm p}_{\mathrm{cf}})^{-1}$缩放真实输入的标签概率并重新归一化。例如三类“时间安排、报名、其他”的$\hat{\bm p}_{\mathrm{cf}}=(0.5,0.3,0.2)$时，原始概率$(0.45,0.35,0.20)$校准为约$(0.293,0.380,0.326)$，预测由“时间安排”改为“报名”。校准只需一次额外前向，但它假设偏置是与输入无关的逐类缩放。

另一条零样本路线把分类改写为自然语言推断：以文档为前提，以“这段文字关于报名”为假设，用推断模型的蕴含概率作为标签分数\scite{nlp-yin2019zeroshot}。它与标签词打分的区别在于模型类型：前者借用为推断而微调的编码器，后者直接使用生成模型的条件概率。立场检测中的标签描述匹配（第~\ref{sec:nlp-classification-stance-nli}小节）是同一思路。

与微调分类器相比，提示式分类不需要训练、能即时接纳新类别；代价是每次调用都要处理说明与示例，大规模在线分类的计算远高于一次BERT前向，结果又对模板、标签词和示例顺序敏感，须在开发集上固定渲染序列并确定阈值。有足够标注且类别稳定时，微调小编码器通常更省也更稳；标注稀少、类别常变或需要先产生银标数据时，提示式分类更合适。

\subsection{多标签与极端多标签分类}
\label{sec:nlp-classification-multilabel}

上一小节已用逐标签sigmoid处理两个标签。标签多起来以后，两个新问题浮现：标签之间有依赖，逐标签独立判断可能给出不合理的组合；标签数达到数万乃至数百万（商品类目、维基百科分类、学科主题词）时，输出层本身的计算与存储成为瓶颈，绝大多数标签又只有寥寥几个正例。

\paragraph{二元相关与标签依赖。}式~\eqref{eq:nlp-classification-sigmoid}把$K$个标签拆成$K$个独立的二分类，称为\term{二元相关}（binary relevance）。它的训练与预测都按标签并行，代价是忽略标签之间的依赖：用逐标签阈值组合出的向量，可能并不是最可能的标签组合。\term{分类器链}（classifier chain）按固定顺序把前面标签的预测作为后面标签的输入，即按链式法则分解联合分布\cite{nlp-read2011classifierchains}：
\[
  p(\bm y\mid x)=\prod_{k=1}^{K}p(y_k\mid x,y_1,\ldots,y_{k-1}).
\]

一个教学数值说明二者的差别。设“线下参加”$y_1$与“线上参加”$y_2$通常二选一：$p(y_1=1\mid x)=0.55$，$p(y_2=1\mid x,y_1=1)=0.15$，$p(y_2=1\mid x,y_1=0)=0.98$。四种组合的联合概率为$(1,0):0.4675$、$(0,1):0.441$、$(1,1):0.0825$、$(0,0):0.009$。二元相关只看边缘概率$p(y_1=1)=0.55$与$p(y_2=1)=0.55\times0.15+0.45\times0.98=0.5235$，阈值0.5下输出$(1,1)$，它的联合概率只有0.0825；链式贪心先取$y_1=1$，再在条件下取$y_2=0$，得到联合众数$(1,0)$。代价是链的顺序须预先选定，前面的错误会传给后面，训练时用真实标签、预测时用预测标签也带来不一致。标签数很大时，链式方法难以扩展，极端多标签方法转而利用标签之间的层次或聚类结构。

\paragraph{极端多标签的挑战。}\term{极端多标签分类}（extreme multi-label classification, XMC）指标签数$L$达到$10^4$到$10^6$以上的多标签任务。设文档表示为$d$维，输出层权重为$L\times d$：$L=670{,}091$（一个商品数据集的标签数）、$d=1536$时约有$1.03\times10^9$个参数，每个样本的前向就要计算全部$L$个得分。另一个困难是长尾：大多数标签的正例只有几个，单独训练的逐标签分类器几乎学不到东西。

\paragraph{XML-CNN。}XML-CNN在TextCNN上做了三处修改，分别针对表示太粗、输出层太大和多标签损失\scite{nlp-liu2017xmlcnn}。第一，\term{动态最大池化}（dynamic max pooling）把每个滤波器的响应序列等分为$p$段，每段各取最大值，一个滤波器输出$p$个数而非一个，保留“关键词出现在文档前部还是后部”的粗粒度位置信息，长文档的多个主题不必挤进同一个最大值。第二，在池化层和输出层之间插入一个较窄的隐藏层（瓶颈层），先把池化特征压到$h$维再映射到$L$个标签。第三，用逐标签二元交叉熵代替softmax，使多个标签可以同时取高概率。

瓶颈层的作用可以直接算参数。取3种宽度、每种128个滤波器、$p=4$，池化特征为1536维；直接接输出层需要$1536L\approx1.03\times10^9$个参数。插入$h=512$的瓶颈后，参数为$1536\times512+512L\approx3.44\times10^8$，降到约三分之一，且瓶颈层由所有标签共享，长尾标签可以借用常见标签学到的表示。代价是所有标签必须通过同一个512维表示，标签再多时这一层仍是信息瓶颈，输出层也仍需对全部标签打分。

\paragraph{AttentionXML。}XML-CNN对所有标签使用同一个文档向量，可是与“机器学习”相关的段落和与“活动场地”相关的段落未必相同。AttentionXML为每个标签计算专属的注意力，并用标签树避免对全部标签打分\scite{nlp-you2019attentionxml}。

模型部分用双向LSTM编码词序列得$\hat{\bm h}_1,\ldots,\hat{\bm h}_T\in\mathbb R^{2N}$，第$j$个标签有自己的注意力参数$\hat{\bm w}_j\in\mathbb R^{2N}$：
\begin{equation}
  \label{eq:nlp-classification-label-attention}
  \alpha_{t,j}=\frac{\exp(\hat{\bm h}_t^\top\hat{\bm w}_j)}{\sum_{t'=1}^{T}\exp(\hat{\bm h}_{t'}^\top\hat{\bm w}_j)},\qquad
  \hat{\bm m}_j=\sum_{t=1}^{T}\alpha_{t,j}\hat{\bm h}_t.
\end{equation}
每个标签得到一个以它为查询的文档表示$\hat{\bm m}_j$，再经所有标签共享的全连接层和输出层得到该标签的概率，损失仍为逐标签二元交叉熵。与HAN的单一查询相比，这里有$L$个查询，同一文档对不同标签呈现不同的“重点”；共享全连接层则控制参数量。

$L$个注意力的时间与存储都与$L$成正比，标签树负责把每个样本需要计算的标签数降下来。AttentionXML使用\term{概率标签树}（probabilistic label tree, PLT）：每个叶子是一个原始标签，内部节点$n$对应一组相似标签，二元变量$z_n$表示“该子树中至少有一个相关标签”。叶子标签的概率按链式法则写成路径上条件概率之积：
\[
  P(z_n=1\mid x)=\prod_{i\in\operatorname{Path}(n)}P\bigl(z_i=1\mid z_{\operatorname{Pa}(i)}=1,x\bigr).
\]
建树时，每个标签用“标注了它的训练文本的词袋特征之和”归一化表示，自顶向下以平衡$k$均值（$k=2$）递归二分，直到每簇标签数不超过$M$；再把这棵深而窄的树压缩成高度为$H$、每个内部节点不超过$K$个子节点的浅而宽的树。训练按层进行，每层一个AttentionXML模型：第一层直接对根的所有子节点训练；第$d$层只对候选节点训练，候选是上一层得分最高的$C$个节点的子节点，因此每层只需$C\times K$个节点的注意力。预测时按同样方式逐层保留前$C$个节点，相当于宽为$C$的束搜索。

以$L=4096$、$K=8$、$H=3$、$C=4$为例：第一层有$4096/8^3=8$个节点，此后每层只计算$4\times8=32$个候选，一个样本共算$8+32\times3=104$个节点，而不是4096个标签。若某叶子路径上的四个条件概率为$0.9,0.8,0.7,0.6$，它的得分为$0.3024$。树的作用与第~\ref{sec:nlp-classical-hierarchical-softmax}小节的层次softmax相似，都是用树结构把对$L$个输出的计算降到近似对数规模；区别在于层次softmax的叶子互斥、概率加和为1，概率标签树的各叶子可以同时为真，每个节点各做一次二分类。代价有两处：层间的错误不可回退，第一层漏掉的子树其下所有标签都不会出现；训练需要$H+1$个模型，原文还用三棵树集成来降低建树随机性的影响。

\paragraph{评价指标。}极端多标签任务通常只关心排在前面的少数标签，评价改用排序指标。设$\bm y\in\{0,1\}^L$为参考标签，$\operatorname{rank}_k(\hat{\bm y})$为得分最高的$k$个标签，
\begin{equation}
  \label{eq:nlp-classification-xml-metrics}
  \begin{aligned}
    \mathrm{P@}k&=\frac1k\sum_{l\in\operatorname{rank}_k(\hat{\bm y})}y_l,\qquad
    \mathrm{DCG@}k=\sum_{r=1}^{k}\frac{y_{(r)}}{\log_2(r+1)},\\
    \mathrm{nDCG@}k&=\mathrm{DCG@}k\Big/\sum_{r=1}^{\min(k,\lVert\bm y\rVert_0)}\frac{1}{\log_2(r+1)},
  \end{aligned}
\end{equation}
$y_{(r)}$是排第$r$位的标签是否相关。nDCG@$k$的分母是理想排序下的DCG，使不同文档之间可比，并对排在前面的命中给更高权重。设参考标签为$\{a,c,e\}$，模型排序为$a,b,c,d,e$：P@1$=1$，P@3$=2/3$，P@5$=3/5$；DCG@3$=1+1/\log_24=1.5$，理想值为$1+1/\log_23+1/\log_24\approx2.131$，nDCG@3$\approx0.704$。这两个指标都偏向常见标签，长尾表现可以另用按标签倾向性加权的版本单独报告。

\subsection{情感与情绪}
\label{sec:nlp-classification-sentiment}

内容分类说明文本在谈什么，评论分析还要说明作者怎样评价所谈对象。\term{情感极性}（sentiment polarity）表示对指定对象的正面、负面或中性评价；\term{情绪}（emotion）表示文本表达的愤怒、喜悦、担忧等状态，可以并存，也可以带强度，二者都不能由正负极性唯一恢复。“服务不错，但价格太高”包含两个方向相反的评价，整篇一个“中性”标签只是相互抵消。

情感分析的挑战有四项：否定与程度修饰改变极性和强度；同一句话对不同目标的评价不同；方面可能是隐含的（“太贵了”评价价格，却没写“价格”）；观点词需要与它评价的方面正确配对。方法依次为：情感词典加否定规则，可解释但覆盖有限；目标条件分类，把评价对象作为输入；方面级情感分析，从给定方面判极性的注意力模型走到一并抽出方面、观点与极性的统一模型。标注还应说明意见持有者，区分作者意见与引述意见；SemEval-2017 Task 4分别设置整条推文情感、目标条件情感与语料级分布估计，说明不同输入粒度要分别定义\scite{nlp-rosenthal2017sentiment}。

\subsubsection{情感词典与否定规则}
\label{sec:nlp-classification-sentiment-lexicon}

可检查的词典基线由极性词得分、程度权重和否定范围构成。对目标$a$的意见词集合$O_a$，一种教学规则是
\[
  S(a)=\sum_{j\in O_a}v_jm_jr_j,
\]
$v_j$是领域词典极性，$m_j$是程度权重，$r_j$是由句法或局部规则确定的否定修正。转折前后、引号内外和不同目标先分组，再汇总。

对“价格不算高，但服务很差”，在价格语境中取“高”的$v=-1$，“不算”修正为$r=-1,m=0.5$，价格得分$+0.5$；“差”取$v=-2$，“很”的$m=1.5$，服务得分$-3$。价格略正面、服务负面；整句求和为$-2.5$，只在业务确实需要整体评价时使用。规则保留了命中词、修正范围和目标，错误容易追查；它的局限也在这里：“服务真好，等了两个小时也没人理”的反讽、“并非不高”的多重否定、“高”在价格与性能中的相反极性，都需要再写规则。

\subsubsection{目标条件情感分类}
\label{sec:nlp-classification-sentiment-bert}

词典规则要人工为每个目标划分意见词。目标条件分类把“评价谁”直接放进模型输入，例如\code{[CLS]}目标\code{[SEP]}评论\code{[SEP]}，训练样本是$(x,a,y)$，同一评论可对应多条目标不同的样本，分类头按极性标签作softmax。

对“服务不错，但价格太高”，标签次序为正面、负面、中性。以“服务”为目标时，教学logits为$(2,0,-1)$，概率约$(0.844,0.114,0.042)$；目标换成“价格”，logits为$(0,2,-1)$，正负两项交换，两条样本的交叉熵都约为0.170。固定文本、只换目标，是检查目标是否真正进入判定的最小测试；若输出不变，模型只利用了整体语气。指定目标不在原文中出现时（“这次安排很周到”对地点并无评价），应允许输出“无评价”，而不是强制选一个极性。

\subsubsection{方面级情感分析：从注意力模型到统一抽取}
\label{sec:nlp-classification-aspect-opinion}

\term{方面级情感分析}（aspect-based sentiment analysis, ABSA）把评论拆成若干方面（服务、价格、场地）及其极性。它包含一组子任务：抽取方面词、抽取观点词、给定方面判极性、抽取方面—观点对，以及一次抽出\term{方面情感三元组}（aspect sentiment triplet）“方面—观点—极性”。下面四类方法依次回答前一类留下的问题。

\paragraph{注意力模型：ATAE-LSTM与IAN。}在LSTM上直接做句子分类，同一句话对“服务”和“价格”只有一个表示。ATAE-LSTM让方面参与两处计算\cite{nlp-wang2016atae}：输入层把方面嵌入$\bm v_a\in\mathbb R^{d_a}$拼接到每个词向量后；输出层以方面为条件做注意力。设LSTM隐状态矩阵$\bm H\in\mathbb R^{d\times N}$，$\bm e_N$为全1向量，
\[
  \begin{aligned}
    \bm M&=\tanh\!\begin{bmatrix}\bm W_h\bm H\\ \bm W_v\bm v_a\bm e_N^\top\end{bmatrix},\qquad
    \bm\alpha=\operatorname{softmax}(\bm w^\top\bm M),\\
    \bm r&=\bm H\bm\alpha^\top,\qquad
    \bm h^*=\tanh(\bm W_p\bm r+\bm W_x\bm h_N),
  \end{aligned}
\]
$\bm h^*$接softmax输出极性。同一句话换一个方面，$\bm\alpha$随之改变，“不错”与“太高”的权重此消彼长。ATAE-LSTM里方面只是一个向量，多词方面（“服务态度”）被平均成一个嵌入，方面内部哪个词重要无从区分。IAN让方面与上下文各用一个LSTM编码并相互注意\cite{nlp-ma2017ian}：以方面隐状态的平均$\bar{\bm t}$为查询，对上下文词打分$\gamma_i=\tanh(\bm h^{c\top}_i\bm W_a\bar{\bm t}+b_a)$，归一化后加权得$\bm c_r$；反过来以上下文平均$\bar{\bm c}$为查询对方面词注意得$\bm t_r$，拼接$[\bm t_r;\bm c_r]$分类。

\paragraph{辅助句构造：BERT-pair。}注意力模型从零学习方面与上下文的交互，训练数据又小。Sun等把方面写成一个辅助句，与评论组成句对输入BERT，用句对分类头判断，借用预训练中学到的句间交互\scite{nlp-sun2019bertpair}。辅助句有四种写法：问句式“你觉得服务怎么样？”；推断式只写方面“服务”；以及两种二元式，对每个候选极性各造一句“服务的评价是正面的”，判断是否成立，一个方面、三种极性就是三条二分类样本。这一步把方面级分类改写为句对分类，与第~\ref{sec:nlp-classification-stance-nli}小节立场的标签描述匹配是同一思路。它和注意力模型一样要求方面事先给定，不能回答“评论里提到了哪些方面”。

\paragraph{统一抽取：BARTABSA、Span-ASTE与GTS。}实际评论的方面和观点都需要从原文中找出来。流水线先抽方面、再抽观点、最后配对判极性，前一步的错误会传给后一步，各子任务还要各自建模。三种统一方法对应第~\ref{sec:nlp-analysis-paradigms}节的三种结构预测范式。

BARTABSA采用约束生成范式，把所有子任务都写成“指针索引与极性类别”组成的序列，用一个带指针的BART编码器—解码器生成\scite{nlp-yan2021bartabsa}。将“服务／不错／但／价格／太高”编号为1至5，三元组任务的目标序列是“1 1 2 2 POS 4 4 5 5 NEG”，每五个记号依次为方面起止、观点起止和极性。解码器每一步的输出空间是“原文位置$\cup$极性类别”：位置的分数由解码器状态与编码器各位置表示的内积给出，类别的分数由解码器状态与类别嵌入的内积给出，这是式~\eqref{eq:nlp-analysis-copy}复制机制只指向原文位置的特例。改变目标序列的格式，同一模型就能完成抽方面、方面—极性、三元组等七个子任务。

Span-ASTE采用跨度范式，直接在方面跨度与观点跨度之间建模，避免逐词元标签无法表达多词方面与多词观点之间的整体配对\scite{nlp-xu2021spanaste}。它枚举宽度不超过$W$的所有跨度，用式~\eqref{eq:nlp-analysis-span-representation}一类的端点拼接加宽度嵌入得到跨度表示，先把每个跨度分为方面、观点或无效；再按两个通道分别保留方面得分与观点得分最高的若干跨度，最后对每个“方面跨度—观点跨度”对拼接两者表示与距离嵌入，分为正面、负面、中性或无关系。剪枝在这里决定可行性：5个词、宽度不超过2时有9个跨度，全部配对为81对；每个通道只留3个，就只需判9对。

GTS采用网格标注，把三元组写进一张词对表\scite{nlp-wu2020gts}。对$i\leq j$的上三角词对$(i,j)$标注六种标签之一：A表示两词属于同一方面，O表示属于同一观点，Pos、Neg、Neu表示一词在方面中、另一词在观点中且二者构成该极性的配对，N表示无关系；对角格$(i,i)$标出单词本身属于方面还是观点。上面的例子中，$(1,1)=\mathrm A$、$(2,2)=\mathrm O$、$(4,4)=\mathrm A$、$(5,5)=\mathrm O$、$(1,2)=\mathrm{Pos}$、$(4,5)=\mathrm{Neg}$，其余为N。解码时先由对角线上连续的A与O恢复方面和观点跨度，再对每个方面—观点对取它们之间各格中出现最多的极性标签。词对表示为两词编码的拼接，原文还让每格的预测迭代读取同行其他格的预测，使方面判断与配对判断互相提示。这种表格思路与第~\ref{sec:nlp-extraction-tplinker}小节关系抽取的TPLinker相同，都把配对关系放到词元对上。

表~\ref{tab:nlp-classification-absa}按子任务覆盖比较这些方法。注意力模型解决“同一句话对不同方面只有一个表示”，辅助句借用预训练交互缓解训练数据小，统一抽取解决“方面须事先给定”和流水线误差传播；三种统一方法分别要付出解析合法序列、枚举跨度对和$O(n^2)$词对表的代价。评价时应分别报告方面与观点的边界F1和完整三元组F1，给定方面后的极性准确率不能代替端到端质量。

\begin{table*}[htbp]
  \centering
  \small
  \begin{tabularx}{\linewidth}{@{}>{\raggedright\arraybackslash}p{24mm}
    >{\raggedright\arraybackslash}p{24mm}>{\raggedright\arraybackslash}X
    >{\raggedright\arraybackslash}X@{}}
    \toprule
    方法 & 输入是否需给定方面 & 输出与核心机制 & 解决的问题与新代价 \\
    \midrule
    ATAE-LSTM、IAN & 需要 & 方面极性；方面条件注意力 & 区分同句不同方面；方面须预先给定，从零学交互 \\
    BERT-pair & 需要 & 方面极性；辅助句与句对分类 & 借用预训练句间交互；仍不抽方面，二元式样本数乘以极性数 \\
    BARTABSA & 不需要 & 全部子任务；指针索引序列生成 & 一个模型覆盖七个子任务；需解析非法序列，解码逐步进行 \\
    Span-ASTE & 不需要 & 三元组；跨度枚举、双通道剪枝、跨度对分类 & 多词方面与观点整体配对；候选数随最大宽度增长 \\
    GTS & 不需要 & 三元组；上三角词对网格标注 & 一张表同时表示跨度与配对；$O(n^2)$词对，边界依赖对角连续性 \\
    \bottomrule
  \end{tabularx}
  \caption{方面级情感分析方法的子任务覆盖。前两行假定方面已知；后三行同时找出方面、观点与极性，分别对应约束生成、跨度和词对表三种结构预测方式。}
  \label{tab:nlp-classification-absa}
\end{table*}

\subsection{目标立场}
\label{sec:nlp-classification-stance}

“改期真让人头疼，不过为了让校外来宾赶上，我支持推迟一天。”语气消极，却对“推迟一天”明确赞成。\term{立场检测}（stance detection）判断文本对指定对象或主张的支持、反对或其他约定关系，输入必须包括目标。情感回答“评价感觉怎样”，立场回答“对这个主张站在哪一边”，证据不同。

立场检测的挑战有三项：目标可能不在文本中出现（评论只谈“来宾赶不上”，却针对“推迟一天”表态）；立场常以隐含方式表达，需要推理才能对上目标；实际应用中的目标不断变化，训练时未见过的新目标没有标注。谱系相应地分三步：目标条件分类直接学习“目标—文本”到立场的映射，需要每个目标都有训练数据；自然语言推断式匹配把立场写成假设，借用推断模型的迁移能力；未见目标方法则专门学习可以跨目标迁移的表示。Bi-STANCE把中英文、名词短语与主张两类目标、未见目标及跨域条件区分开，可用于分别验收这些设定\scite{nlp-zhao2025bistance}。

\subsubsection{BERT目标条件立场分类}
\label{sec:nlp-classification-stance-bert}

将目标和文本组成有序句对，分类头输出支持、反对、中立、无关等任务定义的类别。对上面的评论，以“应推迟一天”为目标，教学logits为$(2,0,-1,-2)$，支持概率约0.831；目标换成“应维持原定日期”，logits变为$(0,2,-1,-2)$，反对概率约0.831。softmax只在同一目标的四个标签内归一化。

这种做法与目标条件情感分类结构相同，难点在数据：目标名称可能泄露训练集的标签偏好（某目标的样本大多是反对），模型会学会只看目标。按目标划分训练与测试，并用目标替换和引述观点（“有人支持推迟，但我没有表态”）的成组样本检查，才能看出模型是否读了文本。

\subsubsection{自然语言推断式标签描述匹配}
\label{sec:nlp-classification-stance-nli}

类别数少而目标常变化时，可以把标签改写成假设：“作者支持将读书会推迟一天”“作者反对将读书会推迟一天”。原文作为前提，用自然语言推断模型（第~\ref{sec:nlp-classification-nli}小节）取得每条假设的蕴含概率。这与第~\ref{sec:nlp-classification-llm}小节的推断式零样本分类是同一机制，模板把固定标签变成可输入的文字。

设支持、反对、中立三条假设的蕴含概率为$(0.75,0.08,0.12)$，它们来自三次独立的推断，不必相加为1。在开发集上确定接受门槛0.6与领先差0.2：最大值0.75超过门槛且领先0.63，输出“支持”；得分为$(0.31,0.29,0.28)$时虽有最大项，却应输出未定。模板措辞、目标中的否定和领域变化都会改变蕴含分数，仍需在立场开发集上校验。

\subsubsection{未见目标的立场检测}
\label{sec:nlp-classification-stance-zeroshot}

推断式匹配依赖通用推断模型，并没有利用已有的立场标注。未见目标立场检测要求在一组目标上训练、在训练中从未出现的目标上测试，问题是怎样让已标注目标的知识迁移到新目标。两种思路分别从“借相近目标”和“去掉目标特有信息”入手。

TGA-Net的思路是：新目标虽未见过，训练集中常有语义相近的目标，可以借用它们\scite{nlp-allaway2020vast}。作者为此构建了VAST数据集：评论来自新闻评论区，目标是数千个不同的名词短语，标签为支持、反对、中立；测试集分零样本目标（训练中未出现）与少样本目标（训练中出现过少量样本）两种设定。模型分三步。第一，把目标与文档作为句对联合输入BERT，得到条件化的目标词元表示$\bm t^{(1)},\ldots,\bm t^{(m)}$与文档词元表示。第二，计算\emph{广义主题表示}：另把文档和目标单独编码并取加权平均得$\bm v_d$、$\bm v_t$，对训练集所有样本的$[\bm v_d;\bm v_t]$做层次聚类；每个样本取最近簇的中心$\bm r_{dt}\in\mathbb R^{2E}$，代表“与它相近的一类目标与文档”。第三，以$\bm r_{dt}$为查询，对目标词元做缩放点积注意力：
\[
  s_i=\operatorname{softmax}_i\bigl(\lambda\,\bm t^{(i)\top}\bm W_a\bm r_{dt}\bigr),\qquad
  \bm c_{dt}=\sum_i s_i\bm t^{(i)},\qquad \lambda=1/\sqrt E,
\]
再把文档词元平均$\tilde{\bm d}$与$\bm c_{dt}$拼接，经两层前馈网络输出三类概率。新目标落入哪个簇，就借用该簇训练样本学到的注意力模式。代价是聚类在训练集上一次完成，簇数需要在开发集上选择，与所有训练簇都不相近的目标得不到帮助。

TOAD从另一侧入手：既然新目标的知识无法预先获得，就让立场表示尽量不含目标特有的信息\scite{nlp-allaway2021toad}。它用BiLSTM以目标为条件编码推文，得到文档向量后经一个线性变换层，再分别送入立场分类器和目标判别器；判别器试图从变换后的表示中认出样本属于哪个目标，训练通过\term{梯度反转}（gradient reversal）让编码器与变换层最大化判别器的损失，同时最小化立场损失。学习目标为
\[
  \mathcal L=\mathcal L_{\mathrm{stance}}-\rho\,\mathcal L_{\mathrm{topic}}
  +\lambda_{\mathrm{rec}}\mathcal L_{\mathrm{rec}}+\lambda_{\mathrm{tr}}\mathcal L_{\mathrm{tr}},
\]
其中$\rho$随训练逐渐增大，重构损失$\mathcal L_{\mathrm{rec}}$防止对抗把编码器的信息全部抹掉，$\mathcal L_{\mathrm{tr}}$把变换层约束在恒等映射附近。判别器还使用新目标的无标注推文，这与领域自适应中“源域有标注、目标域只有无标注数据”的设定相同。评价在SemEval-2016 Task 6的推特数据上逐个留出目标：用其余五个目标训练，在留出目标上测试。

TGA-Net假设新目标与某些训练目标相近，适合目标多、覆盖广的数据；TOAD假设立场表达可以与目标无关地迁移，并需要新目标的无标注文本。二者都难以处理需要背景知识的隐含立场；以目标描述为提示的大语言模型零样本判别是另一种选择，同样须在留出目标上验收。

\subsection{文本来源与生成参与判定}
\label{sec:nlp-classification-origin}

一份流畅的通知可能由人工撰写、模型生成，也可能先由人工写草稿再让模型润色。\term{生成参与判定}（generation involvement detection）估计文本形成过程中模型参与的类别；作者归属则在规定的候选作者之间判断来源，二者不是同一任务。Li等的创作者—编辑者研究区分纯人工、纯模型、模型润色人工文本和经人工化编辑的模型文本\cite{nlp-li2026creator}；实际数据还可能有多轮交替编辑，应按可获得的过程记录标注，必要时增加“混合但过程未知”。

这一任务的挑战比内容分类更尖锐：大模型生成的文本在词汇和句法统计上与人类文本接近；改写可以抹去表面统计特征；检测器换一个生成模型或换一个领域，准确率常明显下降；统计判断也不能当作身份或违规的证明。下面的方法按“需要什么”排列：文体监督分类需要各来源的标注样本；困惑度、DetectGPT与Fast-DetectGPT不训练检测器，但要能访问某个语言模型的概率，后两者以越来越低的调用成本比较原文与其邻近替代文本；水印则要求生成方在生成时主动嵌入信号，换来有统计保证的检测。

\subsubsection{文体特征与来源监督分类}
\label{sec:nlp-classification-origin-style}

可复核的基线从字符$n$-gram、功能词比例、标点频率和句长统计形成向量，标准化后训练线性分类器。作者归属样本需要已知作者、采样时间及题材；生成参与样本还要记录生成系统、提示、编辑方式与程度，标签来自来源记录，而不是另一个检测器的猜测。

取感叹号比例、功能词比例与归一化句长三个特征，某通知的向量为$(0.02,0.30,0.50)$。甲、乙两个候选作者的教学logit为$2x_1-x_2+x_3=0.24$与$-x_1+x_2=0.28$，softmax约$(0.490,0.510)$，低于接受阈值0.8，应输出“未能归属”。上下文编码器可以替换手工特征，但更容易利用题材、平台模板或生成系统专属格式这类捷径，按主题、作者或生成系统留出测试能分别暴露它们。

\subsubsection{困惑度与词元概率检测}
\label{sec:nlp-classification-origin-perplexity}

监督分类需要每个生成系统的样本。若能访问某个语言模型的词元概率，可以不训练检测器，只检查文本是否符合该模型的分布。给定词元序列$w_{1:T}$，平均对数概率与困惑度为
\[
  \bar\ell(x)=\frac1T\sum_{t=1}^{T}\log p(w_t\mid w_{<t}),
  \qquad \operatorname{PPL}(x)=\exp(-\bar\ell(x)).
\]
模型给三个词元的概率为$(0.5,0.25,0.125)$时，平均负对数概率为$\log4$，困惑度为4。模型生成的文本倾向于由高概率词元组成，困惑度偏低。统计前须固定模型、词元化、截断和起止符是否计分，不同词元表下的困惑度不能直接比较。

这个信号混入了题材的影响：常见模板和公式化的人工文本也容易预测，罕见专名则抬高困惑度而不改变来源。困惑度只拿原文与“所有文本”的平均水平比较，下面两种方法改为比较原文与它附近的替代文本。

\subsubsection{DetectGPT：扰动与概率曲率信号}
\label{sec:nlp-classification-detectgpt}

DetectGPT的出发点是：模型生成的文本往往落在模型对数概率的局部高点附近，对它做小幅改写，概率几乎总会下降；人写的文本没有这种倾向，改写后概率可升可降\scite{nlp-mitchell2023detectgpt}。于是用扰动模型$q(\widetilde x\mid x)$（原文用T5）遮住若干跨度并补写，得到$m$份扰动文本，再用待检验的来源模型$p$评分。令$\ell(x)=\log p(x)$，$\bar\ell=m^{-1}\sum_j\ell(\widetilde x_j)$，检测量为
\begin{equation}
  \label{eq:nlp-classification-detectgpt}
  d(x)=\frac{\ell(x)-\bar\ell}
  {\sqrt{\frac{1}{m-1}\sum_{j=1}^{m}
  [\ell(\widetilde x_j)-\bar\ell]^2}}.
\end{equation}
分子近似对数概率在$x$处沿扰动方向的负曲率，分母用扰动分数的标准差归一化，使不同文本之间可比。取四份等长教学改写，原文总对数概率为$-10$，扰动为$(-12,-11,-13,-12)$：均值$-12$，样本方差$2/3$，$d=2/\sqrt{2/3}\approx2.449$，在开发集固定的阈值2.0下判为偏向模型生成。分母为零时应拒绝该次估计。

与困惑度相比，DetectGPT用同一篇文本的邻域作参照，题材和常用程度的影响在分子中大体抵消。代价是计算：每篇文本要调用扰动模型生成数十到上百份改写，再逐份送入来源模型评分；它还假设改写大致保持语义，扰动质量差时信号变弱。

\subsubsection{Fast-DetectGPT：条件概率曲率}
\label{sec:nlp-classification-fast-detectgpt}

DetectGPT的成本来自“整段改写”：每份扰动都是一段新文本，必须分别生成、分别评分。Fast-DetectGPT把扰动换成逐位置的替代词元：在原文前缀固定的条件下，每个位置独立地从采样模型$q_\phi(\cdot\mid x_{<j})$抽一个替代词元，组成样本$\widetilde x$\scite{nlp-bao2024fastdetectgpt}。比较的问题也随之改变：原文在每个位置上选的词元，相对于“同一语境下模型可能选的词元”，是否系统地偏向高概率。其\term{条件概率曲率}（conditional probability curvature）定义为
\begin{equation}
  \label{eq:nlp-classification-fast-detectgpt}
  \begin{aligned}
    d(x)&=\frac{\log p_\theta(x\mid x)-\tilde\mu}{\tilde\sigma},\qquad
    \tilde\mu=\mathbb E_{\widetilde x\sim q_\phi(\cdot\mid x)}\bigl[\log p_\theta(\widetilde x\mid x)\bigr],\\
    \tilde\sigma^2&=\mathbb E_{\widetilde x\sim q_\phi(\cdot\mid x)}\bigl[(\log p_\theta(\widetilde x\mid x)-\tilde\mu)^2\bigr],
  \end{aligned}
\end{equation}
其中$\log p_\theta(\widetilde x\mid x)=\sum_j\log p_\theta(\widetilde x_j\mid x_{<j})$，条件始终是原文前缀，打分模型$p_\theta$与采样模型$q_\phi$可以相同，也可以不同。

去掉扰动成本的关键在于条件独立。各位置的替代词元互不依赖，$\log p_\theta(\widetilde x\mid x)$是独立随机变量之和，均值和方差都按位置分解：
\[
  \begin{aligned}
    \tilde\mu&=\sum_j\tilde\mu_j,\qquad
    \tilde\mu_j=\sum_{v\in V}q_\phi(v\mid x_{<j})\log p_\theta(v\mid x_{<j}),\\
    \tilde\sigma^2&=\sum_j\Bigl[\sum_{v\in V}q_\phi(v\mid x_{<j})\log^2p_\theta(v\mid x_{<j})-\tilde\mu_j^2\Bigr],
  \end{aligned}
\]
$\tilde\mu_j$是第$j$个位置替代词元的期望对数概率。所有位置的分布$q_\phi(\cdot\mid x_{<j})$与$p_\theta(\cdot\mid x_{<j})$都由原文的一次前向计算同时得到，求和遍历词表即可精确计算，不需要真的采样。采样模型与打分模型相同时只需一次前向；分子此时等于原文对数似然加上各位置熵之和，正是两个经典零样本基线的组合。

用一个两位置、三候选词元的教学例走通。设$q_\phi=p_\theta$，第1位置分布为$(0.6,0.3,0.1)$，原文选了第一个；第2位置为$(0.5,0.4,0.1)$，原文选了第二个。逐位置计算得$\tilde\mu_1\approx-0.898$、方差约0.315，$\tilde\mu_2\approx-0.943$、方差约0.216，故$\tilde\mu\approx-1.841$，$\tilde\sigma\approx0.729$。原文$\log p=\log0.6+\log0.4\approx-1.427$，$d\approx0.568$。若第2位置也选了最高概率的词元，$\log p\approx-1.204$，$d\approx0.874$。每个位置都选高概率词元的文本曲率更大，更像模型生成。

与DetectGPT相比，Fast-DetectGPT把“整段扰动加逐份评分”换成“逐位置替代加解析期望”，模型调用从上百次降到一两次，原文报告检测速度提高约340倍，在其测试集上检测效果也更好。它仍然继承了这类方法的共同前提：需要一个与真实来源模型足够接近的打分模型，经过大幅改写或人机共同编辑的文本，信号会被削弱。

\subsubsection{生成水印的嵌入与检测}
\label{sec:nlp-classification-watermark}

前面的方法都在事后猜测来源，信号来自模型自然留下的统计痕迹。若生成方愿意配合，可以在生成时主动嵌入一个只有知道密钥才能检出的信号，这就是\term{文本水印}（text watermarking）。Kirchenbauer等的方法在第~\ref{sec:nlp-generation-sampling}小节的采样流程中只插入一步，修改每一步的logits\cite{nlp-kirchenbauer2023watermark}：以前一个词元$w_{t-1}$与密钥的哈希为随机种子，把词表$V$随机划分为大小$\gamma|V|$的\term{绿名单}（green list）$G_t$和其余的红名单，再给绿名单词元的logits加上常数$\delta>0$：
\[
  \tilde p_t(v)\propto\exp\bigl(l_v+\delta\,\mathbb 1[v\in G_t]\bigr),
\]
$l_v$为原logits，归一化在整个词表上进行。若原分布在绿名单上的总概率为$P_G$，加偏置后变为
\[
  \tilde P_G=\frac{e^\delta P_G}{e^\delta P_G+1-P_G}.
\]取$\gamma=0.25$、$\delta=2$，一个绿名单恰好占四分之一概率的位置，绿词元概率从0.25升到约0.711。检测方只需要密钥和词元化器，不需要语言模型：重算每个位置的绿名单，数出文本中绿词元的个数$|s|_G$。没有水印的文本与密钥无关，每个词元落入绿名单的概率为$\gamma$，$T$个词元中绿词元数近似服从二项分布$\mathrm{Bin}(T,\gamma)$，于是用单侧$z$检验：
\begin{equation}
  \label{eq:nlp-classification-watermark-z}
  z=\frac{|s|_G-\gamma T}{\sqrt{T\gamma(1-\gamma)}}.
\end{equation}
完整流程见算法~\ref{alg:nlp-classification-watermark}。

\begin{algorithm}[绿名单水印的生成与检测]
  \label{alg:nlp-classification-watermark}
  \begin{algorithmic}[1]
    \Require 语言模型logits函数$l(\cdot)$，词表$V$，密钥$\kappa$，哈希函数$H$，绿名单比例$\gamma$，偏置$\delta$，检测阈值$z_0$
    \Procedure{Generate}{$c$}\Comment{$c$为提示}
      \For{$t=1,2,\ldots$直到结束符或长度上限}
        \State 以$H(\kappa,w_{t-1})$为种子，把$V$随机划分为绿名单$G_t$（$\gamma|V|$个）与红名单
        \State 计算$\bm l=l(c,w_{<t})$，对$v\in G_t$令$l_v\gets l_v+\delta$
        \State 按原有解码规则（采样或束搜索）从$\operatorname{softmax}(\bm l)$选出$w_t$
      \EndFor
    \EndProcedure
    \Procedure{Detect}{$w_{1:T}$}
      \State $|s|_G\gets0$
      \For{$t=1$到$T$}
        \State 以$H(\kappa,w_{t-1})$重建$G_t$；若$w_t\in G_t$，则$|s|_G\gets|s|_G+1$
      \EndFor
      \State 按式~\eqref{eq:nlp-classification-watermark-z}计算$z$；$z>z_0$时判为含水印
    \EndProcedure
  \end{algorithmic}
\end{algorithm}

取$T=200$、$\gamma=0.25$，无水印时绿词元的期望为50，标准差$\sqrt{37.5}\approx6.12$。一段含水印的文本若有100个绿词元，$z\approx8.16$；一段人写文本有56个，$z\approx0.98$。阈值$z_0=4$对应的单侧误报率约$3\times10^{-5}$，这个误报率来自零假设下的二项分布，不需要人类文本的训练样本，这是水印与前面所有检测器的根本区别。

水印的强度取决于每一步可供偏移的概率质量。低熵位置几乎无法变绿：模型以0.99的概率选红名单词元时，$\delta=2$只把绿名单总概率从0.01抬到约0.069，代码、固定格式和短答案大量由这类位置组成，同样长度下$z$值较低。改写则逐处抹去信号：种子只依赖前一个词元，替换一个词元最多改变它自身和下一个位置的归属。上例若有40处替换，最多影响80个位置；受影响位置原约一半为绿、改写后按比例$\gamma$变绿时，绿词元约从100降到80，$z\approx4.90$，仍可检出；整段复述则使绿词元比例退回$\gamma$。水印只对配合的生成方有效，$\delta$越大越难抹去，也越可能改变生成质量。

表~\ref{tab:nlp-classification-origin-methods}汇总五种来源判定方法的前提、信号与局限。

\begin{table*}[htbp]
  \centering
  \small
  \begin{tabularx}{\linewidth}{@{}>{\raggedright\arraybackslash}p{24mm}
    >{\raggedright\arraybackslash}X>{\raggedright\arraybackslash}X
    >{\raggedright\arraybackslash}X@{}}
    \toprule
    方法 & 前提 & 核心信号与计算 & 主要局限 \\
    \midrule
    文体监督分类 & 各来源的标注样本 & 字符$n$-gram、功能词等特征上的分类器 & 跨题材、跨生成系统易利用捷径 \\
    困惑度 & 可访问某个语言模型 & 平均对数概率，一次前向 & 混入题材与常见程度，模板文本误报 \\
    DetectGPT & 来源模型与扰动模型 & 原文与整段扰动的对数概率差，数十到上百次调用 & 计算昂贵，依赖扰动保持语义 \\
    Fast-DetectGPT & 打分模型（可另有采样模型） & 逐位置替代词元的解析期望与方差，一两次前向 & 打分模型须接近来源，改写削弱信号 \\
    绿名单水印 & 生成方嵌入，检测方持有密钥 & 绿词元计数的$z$检验，无需语言模型 & 低熵文本信号弱，大量改写可抹去 \\
    \bottomrule
  \end{tabularx}
  \caption{文本来源判定方法对照。前四种在生成之后被动检测，阈值依赖开发集；水印在生成时主动嵌入信号，误报率由零假设下的二项分布给出。输出都是统计判断，不构成作者身份的证明。}
  \label{tab:nlp-classification-origin-methods}
\end{table*}

\section{文本匹配与关系判定}
\label{sec:nlp-classification-matching}

单篇分类固定了标签表，句对任务还要固定关系及其方向。以本章通知为参照，“机器学习读书会改到周五”与它很相似，却省略了日期、地点和报名条件：它可以充当简要概括，却不是保留全部条件的等义复述。完整通知能够推出这句概括，反过来不行。查询“周五几点开始？”与片段“14:00开始”几乎没有共同词，后者却比一段详细的读书会介绍更有用。相似、等价、有向蕴含和对需求有用是四种不同关系，监督数据和验收口径都要分别规定。

这些任务共享一组困难。其一，关系可能有方向，交换两段文本会改变答案。其二，日期、否定、数量和角色这些少量词元决定关系，整体相似度却对它们不敏感。其三，让两段文本逐词比较能提高精度，但每个候选都要重新计算，候选库很大时无法承受。其四，没有人工评分时，需要从别的信号学到可比较的表示。其五，跨语言对应和版本比较还会出现一对多、增删和高重合下的关键改动。本节先比较匹配架构，回答第三个困难；随后按关系类型展开相似度、复述、推断、查询相关性以及对应与版本关系，其余困难在各自的任务中处理。

\subsection{匹配架构的谱系}
\label{sec:nlp-classification-matching-architectures}

所有匹配任务都要为文本对$(a,b)$计算一个分数$s(a,b)$，再按任务把分数变成相似度、标签或排序。架构之间的核心差别在于：两段文本在编码时能否互相看见，以及哪些计算可以在查询到来之前完成。图~\ref{fig:nlp-classification-matching-architectures}按这两个问题把模型分为四类。表示型让两侧独立编码成单向量，交互型把两侧联合编码，后期交互保留多个词元向量、把比较推迟到打分阶段，多向量编码则在查询一侧保留少量向量，再由候选决定如何汇聚。

\begin{figure*}[htbp]
  \centering
  \begin{tikzpicture}[
  x=1mm,y=1mm,figure picture,
  mtinp/.style={figure node,minimum height=6.5mm,minimum width=13mm,
    inner xsep=1pt,draw=FigureInput,fill=FigureInputFill,
    font=\figurefont\scriptsize},
  mtenc/.style={figure node,minimum height=6.5mm,minimum width=13mm,
    inner xsep=1pt,draw=FigureModel,fill=FigureModelFill,
    font=\figurefont\scriptsize},
  mtwide/.style={mtenc,minimum width=30mm},
  mtrep/.style={figure node,minimum height=6.5mm,minimum width=13mm,
    inner xsep=1pt,font=\figurefont\scriptsize},
  mtout/.style={figure node,minimum height=7mm,minimum width=30mm,
    inner xsep=1pt,draw=FigureOutput,fill=FigureOutputFill,
    font=\figurefont\scriptsize},
  mtnote/.style={font=\figurefont\scriptsize,text=FigureMuted,align=center},
  mtfl/.style={figure flow}
]
  % 四个面板：左下 x、宽 35、高 64
  \figurePanel{0}{0}{35}{64}{表示型}{FigureInputFill}
  \figurePanel{38}{0}{35}{64}{交互型}{FigureInputFill}
  \figurePanel{76}{0}{35}{64}{后期交互}{FigureInputFill}
  \figurePanel{114}{0}{35}{64}{多向量（Poly）}{FigureInputFill}

  % 共同的输入行
  \foreach \x in {0,38,76,114}{
    \node[mtinp] (qa\x) at (\x+9.5,10) {查询$a$};
    \node[mtinp] (cb\x) at (\x+25.5,10) {候选$b$};
  }

  % 表示型：两塔独立编码，单向量比较
  \node[mtenc] (e1) at (9.5,22) {编码器};
  \node[mtenc] (e2) at (25.5,22) {编码器};
  \node[mtrep] (u) at (9.5,34) {$\bm u\in\mathbb R^{d}$};
  \node[mtrep] (v) at (25.5,34) {$\bm v\in\mathbb R^{d}$};
  \node[mtout] (s1) at (17.5,46) {$s=\cos(\bm u,\bm v)$};
  \draw[mtfl] (qa0)--(e1); \draw[mtfl] (cb0)--(e2);
  \draw[mtfl] (e1)--(u); \draw[mtfl] (e2)--(v);
  \draw[mtfl] (u)--(s1.south -| u); \draw[mtfl] (v)--(s1.south -| v);

  % 交互型：联合编码
  \node[mtwide] (j) at (55.5,22) {联合编码$[a;b]$};
  \node[mtrep,minimum width=30mm] (h) at (55.5,34) {句对表示$\bm h$};
  \node[mtout] (s2) at (55.5,46) {$s=\bm w^{\top}\bm h$};
  \draw[mtfl] (qa38)--(qa38 |- j.south); \draw[mtfl] (cb38)--(cb38 |- j.south);
  \draw[mtfl] (j)--(h); \draw[mtfl] (h)--(s2);

  % 后期交互：保留词元向量
  \node[mtenc] (l1) at (85.5,22) {编码器};
  \node[mtenc] (l2) at (101.5,22) {编码器};
  \node[mtrep] (qi) at (85.5,34) {$\{\bm q_i\}_{i\le n_q}$};
  \node[mtrep] (dj) at (101.5,34) {$\{\bm d_j\}_{j\le n_d}$};
  \node[mtout] (s3) at (93.5,46) {$\sum_i\max_j\bm q_i^{\top}\bm d_j$};
  \draw[mtfl] (qa76)--(l1); \draw[mtfl] (cb76)--(l2);
  \draw[mtfl] (l1)--(qi); \draw[mtfl] (l2)--(dj);
  \draw[mtfl] (qi)--(s3.south -| qi); \draw[mtfl] (dj)--(s3.south -| dj);

  % 多向量：m 个代码向量汇聚查询
  \node[mtenc] (p1) at (123.5,22) {编码器};
  \node[mtenc] (p2) at (139.5,22) {编码器};
  \node[mtrep] (ym) at (123.5,34) {$\bm y^{1},\ldots,\bm y^{m}$};
  \node[mtrep] (yc) at (139.5,34) {$\bm y_{\mathrm{cand}}$};
  \node[mtout] (s4) at (131.5,46) {按$\bm y_{\mathrm{cand}}$汇聚后点积};
  \draw[mtfl] (qa114)--(p1); \draw[mtfl] (cb114)--(p2);
  \draw[mtfl] (p1)--(ym); \draw[mtfl] (p2)--(yc);
  \draw[mtfl] (ym)--(s4.south -| ym); \draw[mtfl] (yc)--(s4.south -| yc);

  % 预计算说明
  \node[mtnote] at (17.5,3) {候选向量可离线计算};
  \node[mtnote] at (55.5,3) {每个句对重新编码};
  \node[mtnote] at (93.5,3) {离线存词元向量};
  \node[mtnote] at (131.5,3) {离线存单向量};
\end{tikzpicture}

  \caption{四类匹配架构的计算位置。表示型两侧各得一个向量，候选向量可离线算好；交互型把两段文本拼接后联合编码，每个文本对都要完整计算一次；后期交互独立编码但保留全部词元向量，打分时逐个查询词元找最匹配的候选词元；多向量编码（Poly-encoder）把查询压成$m$个向量，再按候选向量加权汇聚。从左到右，打分时允许的交互逐渐增多，离线存储与在线代价的分布也随之改变。}
  \label{fig:nlp-classification-matching-architectures}
\end{figure*}

\subsubsection{表示型匹配}
\label{sec:nlp-classification-representation-matching}

网页搜索要在数以亿计的文档中找出与查询语义相近的若干篇。逐对重新计算不可行，最直接的办法是把每段文本独立映射为一个向量，让向量距离代表语义距离。这类做法称为\term{表示型匹配}（representation-based matching）：编码器$f$先分别产生$\bm u=f(a)$和$\bm v=f(b)$，再用余弦或内积打分；两侧各用一个编码器的结构也称\term{双编码器}（bi-encoder）。由于$\bm v$不依赖$a$，候选库可以提前编码并建立向量索引。

DSSM是这一路线的早期代表\scite{rs-h10}。它首先解决词表过大的问题：把每个英文词两端加上边界符，再切成字母三元组，例如\code{good}变成\code{\#go}、\code{goo}、\code{ood}、\code{od\#}。原文用这种\term{词哈希}（word hashing）把约50万词的词表压缩成约3万维三元组计数，未登录词也能得到表示。查询和文档各自经过三层带$\tanh$的全连接网络，形状依次为$30\text{k}\to300\to300\to128$，最后一层输出即语义向量。对一个查询$q$，取被点击文档$d^+$与4篇随机未点击文档组成候选集合$D_q$，训练目标是
\[
  P(d^+\mid q)=\frac{\exp(\gamma\cos(\bm u_q,\bm v_{d^+}))}
  {\sum_{d\in D_q}\exp(\gamma\cos(\bm u_q,\bm v_d))},
\]
最小化$-\log P(d^+\mid q)$。$\gamma$是平滑因子，把取值在$[-1,1]$的余弦放大成足够分开的logit。设一个查询对点击文档的余弦为0.6，对两篇未点击文档为0.2和0.1，取$\gamma=10$，三项未归一化权重为$(e^6,e^2,e^1)$，点击文档的概率约0.9756，损失约0.0247。这个“一正多负的softmax”就是后文对比损失的雏形。

ARC-I保留同样的独立编码思路，把多层感知机换成卷积\scite{nlp-matching-hu2014arc}。每句的词向量序列经过窗口为$k$的一维卷积与最大池化，堆叠两层后得到定长句向量，两句向量拼接后送入多层感知机输出匹配分。卷积能抓住局部词序，例如“改至周五”与“周五改至”得到不同的局部特征；但两句之间的比较仍然只发生在压缩之后。

表示型的优点是在线代价低：查询编码一次，每个候选只需一次$d$维内积。它的缺陷也来自压缩：句中哪个词与对方哪个词对应，在得到$\bm u,\bm v$时已无法再调整。“周五14:00”和“周六14:00”在单向量中可能只差一个很小的分量，打分阶段无从放大。后文的Sentence-BERT、SimCSE与DPR都沿用这一架构，只是把编码器换成预训练Transformer，并改进训练信号。

\subsubsection{交互型匹配}
\label{sec:nlp-classification-interaction-matching}

要让关键词元之间的对应进入判断，两段文本就必须在编码阶段相遇。\term{交互型匹配}（interaction-based matching）先建立两侧词元之间的交互结构，再从中提取匹配特征。ARC-II是与ARC-I同时提出的交互版本：第一层对句$a$的每个窗口$i$与句$b$的每个窗口$j$，把两段窗口的词向量拼接后做卷积，得到一张以$(i,j)$为坐标的二维特征图，之后交替进行二维池化与卷积，最后由多层感知机打分\scite{nlp-matching-hu2014arc}。

MatchPyramid把这一想法写得更直接：先计算\term{匹配矩阵}（matching matrix）$\bm M\in\mathbb R^{n_a\times n_b}$，元素$M_{i,j}$可以是词相同的指示、词向量余弦或内积，再把它当作一张单通道图像，用二维卷积识别其中的匹配模式\scite{nlp-matching-pang2016matchpyramid}。卷积核在矩阵上滑动，对角方向的连续高值对应有序的短语匹配；动态池化把不同长度句对的特征图池化成固定大小，最后接多层感知机。

用一个可手算的例子说明卷积在看什么。查询为“改至、周五”，候选甲为“现、改至、周五”，候选乙为“周五、现、改至”，用词相同指示构造匹配矩阵：
\[
  \bm M_{\text{甲}}=\begin{pmatrix}0&1&0\\0&0&1\end{pmatrix},\qquad
  \bm M_{\text{乙}}=\begin{pmatrix}0&0&1\\1&0&0\end{pmatrix}.
\]
两张矩阵都有两个1，词袋重合相同。取$2\times2$对角卷积核$\bigl(\begin{smallmatrix}1&0\\0&1\end{smallmatrix}\bigr)$，在两个位置上的响应分别是$M_{1,1}+M_{2,2}$和$M_{1,2}+M_{2,3}$：候选甲为$(0,2)$，候选乙为$(0,0)$。最大池化后甲得2、乙得0，有序短语“改至周五”只在甲中被识别。真实模型的卷积核由训练学到，这里人为给定，只演示结构为何能保留次序。

ESIM（第~\ref{sec:nlp-classification-esim}小节）在软对齐之后比较局部差异，BERT\term{交叉编码器}（cross-encoder，第~\ref{sec:nlp-classification-paraphrase-cross}、\ref{sec:nlp-classification-rerank}小节）则让自注意力在所有层中跨句交互，它们都属于这一类。交互型在复述、推断和重排上通常更准确，代价是没有任何候选侧结果可以复用：$N$个候选就要做$N$次联合计算，每次的规模随$n_an_b$或$(n_a+n_b)^2$增长。因此它多用于对少量已召回候选做细判。

\subsubsection{后期交互与多向量编码}
\label{sec:nlp-classification-late-interaction}

表示型在压缩后才比较，交互型在编码时就比较。能否独立编码、却在打分时保留一定的细粒度交互？\term{后期交互}（late interaction）的回答是不把候选压成单向量。ColBERT独立编码两侧，保留每个词元的低维向量，打分时让每个查询词元在候选中找最相似的词元，再把这些最大值相加\scite{nlp-khattab2020colbert}。候选词元向量可以离线计算，交互只剩下$n_qn_d$次低维内积，不必再运行Transformer；代价是索引要为每个候选保存$n_d$个向量。它的打分公式、训练与存储在第~\ref{sec:nlp-classification-colbert}小节展开。

Poly-encoder从相反的一侧折中：候选仍是单向量，存储与表示型相同；查询一侧保留$m$个向量，让候选决定怎样汇聚\scite{nlp-matching-humeau2020polyencoder}。设查询编码器输出$\bm H=(\bm h_1,\ldots,\bm h_n)\in\mathbb R^{n\times d}$，候选编码器输出$\bm y_{\mathrm{cand}}\in\mathbb R^d$。模型学习$m$个\term{代码向量}（code vector）$\bm c_1,\ldots,\bm c_m\in\mathbb R^d$，每个代码向量对查询词元做一次注意力：
\[
  \bm y^{k}=\sum_{t=1}^{n}w_{t}^{(k)}\bm h_t,\qquad
  \bigl(w_{1}^{(k)},\ldots,w_{n}^{(k)}\bigr)=
  \operatorname{softmax}\bigl(\bm c_k^\top\bm h_1,\ldots,\bm c_k^\top\bm h_n\bigr),
\]
得到$m$个查询向量$\bm y^1,\ldots,\bm y^m$。这一步与候选无关，查询到来时只做一次。随后由候选向量对它们再做一次注意力：
\[
  \begin{aligned}
    (a_1,\ldots,a_m)&=\operatorname{softmax}\bigl(\bm y_{\mathrm{cand}}^\top\bm y^1,\ldots,\bm y_{\mathrm{cand}}^\top\bm y^m\bigr),\\
    \bm y_{\mathrm{ctxt}}&=\sum_{k=1}^{m}a_k\bm y^{k},\qquad
    s=\bm y_{\mathrm{ctxt}}^\top\bm y_{\mathrm{cand}}.
  \end{aligned}
\]
每个候选的在线代价是$m$次内积、一次$m$项softmax和一次加权求和，即$O(md)$。

取一个二维教学例子。查询有两个词元向量$\bm h_1=(1,0)$、$\bm h_2=(0,1)$，代码向量为$\bm c_1=(2,0)$、$\bm c_2=(0,2)$。第一组注意力权重为$\operatorname{softmax}(2,0)\approx(0.8808,0.1192)$，故$\bm y^1\approx(0.8808,0.1192)$，同理$\bm y^2\approx(0.1192,0.8808)$。候选$\bm y_{\mathrm{cand}}=(1,0)$与两者的内积为$(0.8808,0.1192)$，softmax权重约$(0.6817,0.3183)$，汇聚得$\bm y_{\mathrm{ctxt}}\approx(0.6383,0.3617)$，分数约0.6383。候选$(0,1)$按对称计算也得0.6383，但所用的查询向量是$(0.3617,0.6383)$。若改用平均池化的单向量$(0.5,0.5)$，两个候选都只得0.5。可见查询表示会随候选偏向与之相关的那部分内容；$m=1$时Poly-encoder退化为表示型，$m$增大则向更细的交互靠近。

表~\ref{tab:nlp-classification-matching-architectures}汇总四类架构。复杂度以$N$个候选、查询长$n_q$、候选长$n_d$、向量维数$d$计，忽略对所有架构相同的查询编码。精度一栏只给文献中常见的相对位置：Poly-encoder原文报告其效果介于双编码器与交叉编码器之间，ColBERT原文报告其在段落重排上接近BERT交叉编码器而延迟低得多；具体差距随数据和训练而变，不构成统一排名。

\begin{table*}[htbp]
  \centering
  \small
  \begin{tabularx}{\linewidth}{@{}>{\raggedright\arraybackslash}p{12mm}
    >{\raggedright\arraybackslash}p{30mm}>{\raggedright\arraybackslash}p{19mm}
    >{\raggedright\arraybackslash}p{15mm}>{\raggedright\arraybackslash}p{30mm}
    >{\raggedright\arraybackslash}p{10mm}>{\raggedright\arraybackslash}X@{}}
    \toprule
    架构 & 代表模型 & 编码时交互 & 候选预计算 & 每候选在线代价 & 存储 & 精度位置与典型用途 \\
    \midrule
    表示型 & DSSM、ARC-I、SBERT、SimCSE、DPR、LaBSE & 否 & 可，单向量 & $O(d)$内积，可用近似近邻 & $Nd$ & 最低一档；大规模召回、去重 \\
    交互型 & ARC-II、MatchPyramid、可分解注意力、ESIM、BERT交叉编码器 & 是，逐层或逐词 & 否 & 完整前向，约$O(n_qn_dd)$至$O(L(n_q+n_d)^2d)$ & 无 & 最高一档；复述、推断、重排 \\
    后期交互 & ColBERT & 否，打分时逐词元比较 & 可，词元向量 & $O(n_qn_dd)$内积，无Transformer & $Nn_dd$ & 接近交叉编码；召回或重排 \\
    多向量 & Poly-encoder & 否，打分时对$m$个向量注意力 & 可，单向量 & $O(md)$ & $Nd$ & 介于表示型与交互型之间；对话回复选择 \\
    \bottomrule
  \end{tabularx}
  \caption{匹配架构的谱系。“编码时交互”一栏决定精度上限，“候选预计算”一栏决定能否面向大候选库。$L$为Transformer层数。实际系统常把它们串联：表示型或后期交互召回少量候选，交互型再对这些候选重排。}
  \label{tab:nlp-classification-matching-architectures}
\end{table*}

\FloatBarrier
\subsection{语义相似度}
\label{sec:nlp-classification-similarity}

\term{语义文本相似度}（semantic textual similarity，STS）给句子对一个连续分数，表示两句在约定标注规范下有多接近。STS基准通常用0到5的人工评分，从“无关”到“完全等义”\scite{nlp-intro-cer2017}。这里有两个困难：相似度对日期、否定等少量关键差异往往不敏感；无标注或标注很少时，怎样得到好的句向量。方法谱系沿此展开：TF-IDF与平均词向量不需训练，给出基线；Sentence-BERT用孪生结构把预训练编码器调成可比较的句向量；SimCSE用对比学习在无标注或NLI数据上进一步改善向量空间。三者都是表示型架构。

\subsubsection{TF-IDF与词向量相似度基线}
\label{sec:nlp-classification-similarity-baseline}

沿用训练语料拟合的词表和IDF，两句表示为$\bm x,\bm y$，非零向量的余弦为
\[
  s(\bm x,\bm y)=\frac{\bm x^\top\bm y}{\|\bm x\|_2\|\bm y\|_2}.
\]
若只保留“读书会、周五、周六”三个等IDF维度，两句向量为$(1,1,0)$与$(1,0,1)$，余弦为$1/2$。若词表把日期词过滤掉，两句都只剩“读书会”，分数反而变成1。空向量需要预先约定缺省分数。平均词向量把词元向量相加后除以词数，也可按IDF加权；“甲通知乙”与“乙通知甲”的平均完全相同，因为求和不保留角色。这两种基线适合快速找同主题候选。

评价时，Pearson相关系数比较预测与参考分数的线性关联，Spearman相关系数比较排序。参考为$(1,2,3)$、预测为$(11,12,13)$时两种系数都为1，分数的量纲却完全错开。若部署要求$s\geq t$才接受，还须在开发集选定$t$，再在测试集报告该阈值下的误报与漏报。

\subsubsection{Sentence-BERT：孪生编码与句向量}
\label{sec:nlp-classification-sbert}

BERT交叉编码器在STS上很准，但为$n$句两两找最相似对需要$n(n-1)/2$次联合前向；直接取BERT的分类位置向量或平均输出做余弦，效果又不好。Sentence-BERT（SBERT）用\term{孪生网络}（siamese network）解决这一矛盾：同一个BERT编码器分别处理两句，对输出作平均池化，得到$\bm u=f_\theta(a)$、$\bm v=f_\theta(b)\in\mathbb R^{d}$，再用句对监督微调，使余弦本身就有意义\scite{nlp-reimers2019sbert}。共享参数不等于拼接输入：每句在编码时只看到自身，因此句向量可以离线计算。

原文给出三种目标。分类目标把$[\bm u;\bm v;|\bm u-\bm v|]\in\mathbb R^{3d}$乘以$\bm W\in\mathbb R^{k\times3d}$后取softmax，在NLI三分类上训练；回归目标以$\cos(\bm u,\bm v)$拟合STS分数，用均方误差；三元组目标对锚点$a$、正例$p$、负例$n$最小化
\[
  L_{\mathrm{triplet}}=\max\{0,\|f(a)-f(p)\|_2-\|f(a)-f(n)\|_2+\epsilon\}.
\]
以回归为例，归一化后的通知向量$\bm u=(1,0)$、近义句$\bm v=(0.8,0.6)$，余弦为0.8；若参考分数缩放到$[0,1]$后为0.9，单对平方误差为0.01。以三元组为例，若负例为$(0.6,0.8)$、$\epsilon=0.5$，正、负距离为$\sqrt{0.4}$和$\sqrt{0.8}$，损失约0.2380：正例虽然更近，但尚未拉开到间隔要求。

三元组每次只比较一个负例。很多数据只有正对，例如问题与答案、标题与正文、原句与译句；这时可以让同一批次里其他样本的正例充当负例，一次比较$B-1$个负候选。设一批有$B$个正对$(q_i,p_i)$，每个$q_i$还可附加\term{难负例}（hard negative）集合$\mathcal H_i$，即与$q_i$相似却不匹配、模型容易混淆的候选，候选集合为$C_i=\{p_1,\ldots,p_B\}\cup\mathcal H_i$。\term{批内负例}（in-batch negatives）对比损失为
\begin{equation}
  L_i=-\log\frac{\exp\bigl(s(q_i,p_i)/\tau\bigr)}
  {\sum_{c\in C_i}\exp\bigl(s(q_i,c)/\tau\bigr)},
  \label{eq:nlp-matching-infonce}
\end{equation}
批损失为$\frac1B\sum_iL_i$。$s$是余弦或内积，$\tau>0$是温度，分母包含正例本身。这一形式常称为InfoNCE损失\cite{nlp-matching-oord2018cpc}：它把匹配写成“在$|C_i|$个候选中认出正例”的分类问题。DSSM的一正四负softmax是它的特例，$\gamma$对应$1/\tau$。它与第~\ref{sec:nlp-classical-nce}小节噪声对比估计的思路相近，都是用与噪声样本的区分代替完整归一化，但这里不估计语言模型的归一化常数，负例来自数据本身而非人为噪声分布。sentence-transformers等实现中常用的“多负例排序损失”就是式~\eqref{eq:nlp-matching-infonce}在余弦与固定缩放下的形式。

对$s(q_i,c)$求导，正例的梯度是$(P_i(p_i)-1)/\tau$，负例$c$的梯度是$P_i(c)/\tau$，其中$P_i(c)$为该候选的softmax概率。因此损失最在意那些已经得分较高的负例。设一个锚点对正例和两个负例的余弦为$(0.8,0.6,0)$。$\tau=0.2$时未归一化权重为$(e^4,e^3,1)$，正例概率约0.7214，损失约0.3266，余弦0.6的负例分得约0.2654的概率，是主要的推开对象。$\tau=1$时正例概率只有约0.4409，两个负例的概率相差不大。温度越低，梯度越集中在难负例上。批量$B$越大，负例越多，估计越接近在全库中区分正例；代价是显存。若批内恰有与$q_i$同义的样本，它会被当成负例推开，这类\term{假负例}（false negative）需要靠去重和采样控制。

\subsubsection{SimCSE：句向量对比学习}
\label{sec:nlp-classification-simcse}

没有人工相似度时，正对从哪里来？无监督SimCSE的回答是同一句话本身：把句子$x_i$送入带dropout的编码器两次，两次随机掩码不同，得到$\bm h_i$与$\bm h_i^+$；以$(\bm h_i,\bm h_i^+)$为正对、批内其他句子的第二次编码为负例，代入式~\eqref{eq:nlp-matching-infonce}，用余弦和$\tau=0.05$训练\scite{nlp-gao2021simcse}。dropout在这里充当最小的数据增强，它没有产生新的同义句，只要求表示对编码扰动稳定、又能与其他句子区分。原文用对齐性（alignment）与均匀性（uniformity）分析效果：正对越近、全体向量在单位球面上分布越均匀，句向量越好用；预训练模型原始句向量挤在狭窄锥体里，对比目标主要改善了均匀性。

有监督SimCSE使用NLI数据：前提为锚点，蕴含句为正例，矛盾句为难负例，即式~\eqref{eq:nlp-matching-infonce}中$\mathcal H_i$只含该锚点的矛盾句，其余锚点的正例与矛盾句也进入分母。训练信号由“同句不同扰动”变成了已标注的意义关系。它仍输出对称的余弦相似度，不会因为用了蕴含句就变成有方向的蕴含判断；运行时关闭dropout，每句只编码一次。检验时仍应单列“甲通知乙／乙通知甲”“须报名／无需报名”一类反例组：平均相关性高，不保证关键条件变化时分数会充分下降。

\subsection{复述与语义等价}
\label{sec:nlp-classification-paraphrase}

\term{复述}（paraphrase）指不同表达在任务规定范围内传达相同含义。按“保留活动事实和适用条件”验收时，“校外来宾须在4月16日18:00前提交报名信息”与“校外来宾提交报名信息的时间须早于4月16日18:00”是正例；删去“校外”会扩大对象，删去时刻会改变截止条件，都应构造为负例。若业务只需要概括，应另定近似复述规范。

复述有两个特有的难点。词面重合高、意义却不同的难例很常见，PAWS用打乱词序构造这类句对，例如“甲从海棠楼到图书馆”与“甲从图书馆到海棠楼”\cite{nlp-intro-zhang2019paws}；等价关系要求对称，交换两句不应改变判定。双编码器速度快，适合候选初筛；交叉编码器让两句逐词交互，能区分角色交换与条件遗漏，用于细判。

\subsubsection{双编码器复述判定}
\label{sec:nlp-classification-paraphrase-bi}

在句向量$\bm u,\bm v$之上构造关系特征并训练二分类头：
\[
  \bm r=[\bm u;\bm v;|\bm u-\bm v|;\bm u\odot\bm v],
  \qquad p=\sigma(\bm w^\top\bm r+b),
\]
$\odot$为逐维乘积，损失为二元交叉熵。增加乘积项是这里的关系头选择，并非SBERT原文的分类配置。前两段若使用不同权重，交换输入会改变结果；令对应权重相等、只用$|\bm u-\bm v|$与$\bm u\odot\bm v$这类对称特征，或平均两个输入次序的预测，都能保证对称。

取对称头：原向量权重为零，差向量权重$(-2,-2)$，乘积权重$(3,3)$，偏置$-0.5$。$\bm u=(1,0)$与保留全部条件的改写$\bm v=(0.96,0.28)$给出$z=-2(0.04+0.28)+3(0.96)-0.5=1.74$，$p\approx0.8507$；删去关键条件的候选$(0.8,0.6)$给出$z=0.3$、$p\approx0.5744$。在开发集选定阈值0.8后，前者接受、后者拒绝。但若角色相反的句子被编码成与$\bm u$相同的$(1,0)$，同一头给出$\sigma(2.5)\approx0.9241$，任何阈值都无法将它与真实等义句分开。这是表示碰撞，只能反馈给表示训练、难负例采样或交给交叉编码器。

\subsubsection{交叉编码器复述判定}
\label{sec:nlp-classification-paraphrase-cross}

交叉编码器把两句作为\code{[CLS]}句A\code{[SEP]}句B\code{[SEP]}联合编码，让“出发地”“目的地”等词元在每层注意力中相互比较，再由\code{[CLS]}表示经线性层预测二元标签，损失仍是二元交叉熵；任务头的一般写法见第~\ref{sec:nlp-classical-bert-heads}小节。它不要求先把每句压成单向量，代价是无法为候选预存表示。

为展示分类头怎样使用交互结果，把一个已训练的三维句对表示暂解释为“角色一致、条件缺失、否定冲突”三个线索，设$z=4h_1-3h_2-5h_3-2$。这只是可手算的教学坐标。完整改写$\bm h=(1,0,0)$给出$p\approx0.8808$，交换出发地和目的地的$(0,0,0)$给出0.1192，遗漏报名条件的$(1,1,0)$给出0.2689，否定冲突的$(1,0,1)$给出0.0474。验证时应成组替换实体、数量、否定和截止条件，统计固定阈值下的结果。联合编码对输入次序敏感，对称任务需双向训练或平均两个次序的分数；长度截断可能切掉第二句末尾的条件，须报告截断率。

\subsection{自然语言推断}
\label{sec:nlp-classification-nli}

\term{自然语言推断}（natural language inference，NLI）给定前提$P$和假设$H$，判断$P$是否支持$H$（蕴含）、与$H$不相容（矛盾），或不足以决定（中立）。三类判断依据数据集约定的通常语义，不是形式证明。以完整通知为前提，“读书会本次在4月18日14:00开始”被蕴含，“读书会本次仍在4月17日开始”与之矛盾，“读书会由张老师主讲”是中立；交换前提与假设，短句推不出报名条件，关系转为中立。MultiNLI扩展体裁，ANLI用对抗式收集增加难例，XNLI考察跨语言迁移，使用时应分别报告训练与评价语言。

NLI有三个困难：关系有方向；否定与量词改变结论，例如“所有校外来宾须报名”推不出“所有参加者须报名”；标注产物会留下捷径，只看假设就能猜出不少标签。方法谱系从轻到重：可分解注意力用软对齐加局部比较，参数少且可并行；ESIM加入BiLSTM上下文编码与推断组合层；BERT类交叉编码直接联合编码，生成式模型以标签词或解释输出。三者都属于交互型架构。第~\ref{chap:nlp-question-answering-dialogue}章会在这一接口之外增加证据检索与主张核验。

\subsubsection{可分解注意力模型}
\label{sec:nlp-classification-decomposable}

可分解注意力模型把句对判断拆成对齐（attend）、比较（compare）和聚合（aggregate）三步\scite{nlp-parikh2016decomposable}。对前提词表示$\bm a_i$与假设词表示$\bm b_j$，前馈网络$F$给出对齐分数，再按行、按列归一化得到软对齐：
\[
  \begin{gathered}
    e_{i,j}=F(\bm a_i)^\top F(\bm b_j),\\
    \widetilde{\bm b}_i=\sum_j\frac{\exp(e_{i,j})}{\sum_k\exp(e_{i,k})}\bm b_j,\qquad
    \widetilde{\bm a}_j=\sum_i\frac{\exp(e_{i,j})}{\sum_k\exp(e_{k,j})}\bm a_i.
  \end{gathered}
\]
比较网络$G$分别处理$[\bm a_i;\widetilde{\bm b}_i]$与$[\bm b_j;\widetilde{\bm a}_j]$，两侧各自求和得$\bm v_1,\bm v_2$，再由$H([\bm v_1;\bm v_2])$三分类。各词对的计算互不依赖，可完全并行；基础版本不含词序，句内注意力是可选增强。

把前提简化为“读书会”$(1,0)$和“14:00”$(0,1)$，假设只有“读书会”$(1,0)$，两条对齐分数为$(\log3,0)^\top$。按行归一化时每个前提词只能对齐唯一的假设词，权重都是1；按列归一化时，假设词对前提的权重为$(3/4,1/4)$，得$\widetilde{\bm a}=(0.75,0.25)$。前提中“14:00”与其对齐对象的绝对差为$(1,1)$，假设一侧的差为$(0.25,0.25)$。若$G$暂取绝对差的$L_1$和，两侧聚合值为$2$和$0.5$：前提有假设没有的内容，反之不然，这种不对称正是蕴含方向的线索。例子也暴露了一个缺陷：单候选时“14:00”被迫对齐到“读书会”，软对齐不会自动拒绝无对应的词；没有词序，角色交换也难以区分。

\subsubsection{ESIM：上下文编码与局部推断组合}
\label{sec:nlp-classification-esim}

ESIM针对上述两点改进\scite{nlp-chen2017esim}。第一层用BiLSTM编码每句，使每个位置的表示包含句内语境；按可分解注意力的方式计算软对齐后，对前提一侧构造增强特征
\[
  \bm m_i=[\bm a_i;\widetilde{\bm b}_i;\bm a_i-\widetilde{\bm b}_i;\bm a_i\odot\widetilde{\bm b}_i]\in\mathbb R^{8h},
\]
$h$为单向隐藏维数，差与积显式给出“哪里一致、哪里冲突”。另一侧同样处理。增强特征经前馈投影后进入第二层BiLSTM，把分散的局部比较组合起来；对两侧输出分别做平均池化与最大池化，四个向量拼接后经多层感知机三分类。BiLSTM的原理可参考\BookChapterRef{chap:recurrent-networks}{卷五的循环神经网络章}，这里只用到它的输出在每个位置同时汇总了左右语境。

比较“甲通知乙报名”与“乙通知甲报名”时，词相同，但第一层BiLSTM使“甲作为通知者”与“甲作为接收者”的表示不同。设某位置编码为$\bm a=(0.8,0.2)$，对齐表示为$\widetilde{\bm b}=(0.3,0.7)$，则差为$(0.5,-0.5)$，积为$(0.24,0.14)$，增强特征为$(0.8,0.2,0.3,0.7,0.5,-0.5,0.24,0.14)$，其中较大的差值标记了角色不一致。若第二层输出两个位置为$(0.2,0.8)$、$(0.6,0.4)$，平均池化为$(0.4,0.6)$，最大池化为$(0.6,0.8)$。与可分解注意力相比，ESIM用顺序计算换来了语境信息；它仍可能学到只看词重合的捷径。

\subsubsection{BERT类模型与生成式关系判定}
\label{sec:nlp-classification-nli-bert}

BERT类交叉编码器把前提与假设按次序拼接联合编码，分类头以三分类交叉熵训练。若三条假设分别得到logits $(3,0,0)$、$(0,3,0)$、$(0,0,3)$，每条首选标签的概率都是$e^3/(e^3+2)\approx0.9094$，损失约0.0949。交换前提和假设需要重新编码，不能直接对调输出概率。受约束生成让模型只输出三个合法标签词之一，监督目标是标签词元的条件似然；自由生成解释是另一项输出，须单独核查它是否引用了前提中的实际内容。诊断时应比较只给假设的模型、词汇重合基线和完整句对模型，前两者若已取得高分，说明数据存在标签泄漏。

\subsection{查询与候选文本的相关性}
\label{sec:nlp-classification-relevance}

检索中的\termalias{search-relevance}{相关性}是候选对当前信息需求的有用程度。设三个教学候选为$d_1$“机器学习读书会介绍”、$d_2$“ML-0418更正：4月18日14:00开始”、$d_3$“原计划4月17日举行”。对查询“读书会几点开始？”，$d_2$是答案来源，$d_1$只与主题相关，$d_3$词面相关却已过时。挑战因此有三个：查询与答案用词不同；候选库很大，必须先召回再精排；相关与时效可能冲突。谱系依次是：BM25以词汇统计高效召回，但无法处理同义；DPR类双编码器学习语义召回，难负例挖掘解决随机负例过易的问题；交叉编码器在小候选集上精排；ColBERT以后期交互兼顾预计算与细粒度匹配。倒排索引、近似近邻检索、SPLADE等稀疏学习表示以及在线系统架构，可参考\BookChapterRef{chap:rs-search}{卷十一的搜索相关性章}；本节只写这些模型在文本匹配中的结构与训练，第~\ref{chap:nlp-question-answering-dialogue}章的检索问答回用这里的模型。

\subsubsection{BM25：词汇相关性基线}
\label{sec:nlp-classification-bm25}

BM25按查询词在候选中的出现、词项区分度与文档长度打分\scite{rs-h01}。采用非负IDF版本：
\[
 \operatorname{BM25}(q,d)=\sum_{t\in q}
 \log\!\left(1+\frac{N-\operatorname{df}_t+0.5}{\operatorname{df}_t+0.5}\right)
 \frac{f_{t,d}(k_1+1)}{f_{t,d}+k_1(1-b+b|d|/\overline L)},
\]
$N$为文档数，$\operatorname{df}_t$为含词$t$的文档数，$f_{t,d}$为词频，$\overline L$为平均长度；$k_1$控制词频饱和，$b$控制长度归一。取$N=4$、$\operatorname{df}=2$、$k_1=1.2$、$b=0.75$、$|d|=\overline L$，该词IDF为$\log2$；出现1次贡献0.6931，出现2次贡献$\log2\times4.4/3.2\approx0.9531$，并未翻倍。BM25能奖励“读书会”的匹配，却不知道“何时开始”与“几点举行”同义，也不知道哪个日期已被更正。中文分词、活动编号是否拆开和时间规范化都会改变匹配，须作为索引配置记录。

\subsubsection{DPR类双编码器与难负例挖掘}
\label{sec:nlp-classification-dpr}

DPR用两个BERT编码器分别编码查询与段落，以分类位置向量的内积$s(q,d)=f_q(q)^\top f_d(d)$打分，训练目标是式~\eqref{eq:nlp-matching-infonce}取$\tau=1$、内积打分的形式\scite{nlp-karpukhin2020dpr}。原文比较了随机负例、BM25高分负例和批内负例，最终把批内负例与每个问题一条BM25难负例结合。设查询向量为$(1,1)$，$d_2,d_1,d_3$的向量为$(1,1),(1,0),(0,0)$，内积为$(2,1,0)$，正例概率$e^2/(e^2+e+1)\approx0.6652$，损失约0.4076。部署时离线编码全部段落并绑定文档编号、字符区间和版本，查询在线编码后取前$k$个，Recall@$k$衡量正确段落是否进入前$k$。召回阶段漏掉的段落，后续重排无法找回。

负例的难度决定模型学到什么。随机负例与查询几乎无关，模型很快就能把它们分开，梯度随之消失；BM25负例与查询词面相近，却未必是当前模型最容易混淆的样本。难负例挖掘的发展沿着“负例是否来自当前模型的错误”展开。ANCE直接用正在训练的模型检索负例：用当前编码器为整个语料建近似近邻索引，对每个查询取排名靠前的非正例段落作负例\scite{rs-s03}。困难在于编码器每更新一步，全库向量就已过时，而重新编码数百万段落代价很高。ANCE的做法是异步刷新：训练进程持续更新参数，另一个推断进程按固定间隔取最新检查点重新编码语料、重建索引，完成后替换旧索引，见算法~\ref{alg:nlp-classification-ance}。训练用的负例因此略为滞后，但始终接近当前模型的错误。

\begin{algorithm}[ANCE式异步难负例训练]
  \label{alg:nlp-classification-ance}
  \begin{algorithmic}[1]
    \Require 训练对$\{(q,d^+)\}$，语料$\mathcal C$，以BM25负例预训练的编码器参数$\theta_0$，每查询负例数$k$，刷新间隔$m$
    \Ensure 编码器参数$\theta_T$
    \State 用$\theta_0$编码$\mathcal C$，建立近似近邻索引$\mathcal I$
    \For{$t=1,\ldots,T$}
      \State 取一批查询；对每个$q$，从$\mathcal I$检索前若干段落，去掉已标注正例，抽取$k$个作为$\mathcal H_q$
      \State 以式~\eqref{eq:nlp-matching-infonce}（候选为批内正例与$\mathcal H_q$）计算损失，更新$\theta_t$
      \If{$t\bmod m=0$}\Comment{与训练并行执行}
        \State 推断进程用检查点$\theta_t$重新编码$\mathcal C$，建好新索引后替换$\mathcal I$
      \EndIf
    \EndFor
  \end{algorithmic}
\end{algorithm}

难负例越难，混入假负例的可能越大：排名最靠前的“负例”可能恰好是未被标注的正确段落。另一类方法借助交叉编码器的判断。RocketQA用交叉编码器给检索出的难负例重新打分，剔除得分很高的疑似假负例，并用它为未标注问题生成伪标签\scite{nlp-matching-qu2021rocketqa}。\term{知识蒸馏}（knowledge distillation）则直接让双编码器学生模仿交叉编码器教师的打分。Margin-MSE只要求学生复现教师对正负候选的分差：
\[
  L_{\mathrm{MM}}=\Bigl(\bigl[s_S(q,d^+)-s_S(q,d^-)\bigr]-\bigl[s_T(q,d^+)-s_T(q,d^-)\bigr]\Bigr)^2,
\]
$s_S,s_T$分别为学生与教师分数\scite{nlp-matching-hofstatter2020margin}。只匹配分差，是因为两种架构的分数尺度不同，绝对值无法直接对齐。若教师分差为3、学生为1，损失为4，梯度推动学生拉大分差。蒸馏把交互型的精度部分转移到表示型中，推断时仍只需内积；它的上限受教师质量限制，教师的系统偏差也会一并传给学生。

\subsubsection{BERT交叉编码器：候选重排}
\label{sec:nlp-classification-rerank}

重排器只处理已召回的候选，把查询与每个段落拼接后联合编码，由\code{[CLS]}表示输出相关分\scite{rs-s20}。点式训练用二元交叉熵；成对训练最小化$\log(1+\exp[-(s^+-s^-)])$，要求正例高于负例；列表式训练对一组候选用softmax交叉熵。设初始列表为$(d_1,d_3,d_2)$，重排器给出logits $(0,-1,2)$，排序恢复为$(d_2,d_1,d_3)$。以“能回答当前开始时间”为唯一相关标准，正确段落的倒数排名从$1/3$升至1；MRR是跨查询的平均倒数排名，分级相关性下则用nDCG并注明截断$k$。时间难负例要保留“原定”“现改至”等关系词，而不只是随机换日期。若召回集合里只有$d_1,d_3$，重排器仍会排出第一名，因此还需要“无合适候选”的拒绝规则。

\subsubsection{ColBERT：词元表示与后期交互}
\label{sec:nlp-classification-colbert}

单向量可能压掉稀少而关键的词，交叉编码又太贵。ColBERT用共享BERT分别编码查询与文档，在开头分别加查询、文档标记；查询用掩码词元补到固定长度，让模型学会查询扩展，文档一侧滤掉标点的向量；输出经线性层降到较低维度（原文128维）并作$L_2$归一化\scite{nlp-khattab2020colbert}。打分采用MaxSim聚合：
\[
  S(q,d)=\sum_{i=1}^{n_q}\max_{1\le j\le n_d}\bm q_i^\top\bm d_j.
\]
每个查询词元找最匹配的文档词元，同一个文档词元可以被多次选中。给定相似矩阵
$\bm M=\bigl(\begin{smallmatrix}0.8&0.2&0.1\\0.1&0.3&0.9\end{smallmatrix}\bigr)$，
行最大值为0.8与0.9，MaxSim为1.7；若改为平均池化后的单向量内积，结果等于矩阵全部元素的均值0.4，每个查询词元的最佳对应被平均掉了。

训练用查询、正段落和负段落三元组上的成对softmax交叉熵。文档词元向量离线计算并存入索引，查询到来后只做矩阵乘法与按行取最大，因此既能用于重排，也能结合近似近邻检索直接召回。代价是存储：每篇文档存$n_d$个向量，而非一个，ColBERTv2等后续工作用残差压缩降低这一开销。局部最大匹配也有盲区：“原时间14:00，新时间15:00”中两个时间词都可能得到高分，各词元高分不等于组成了一致答案。

\subsection{文本对应、重复与版本关系}
\label{sec:nlp-classification-alignment-duplicates}

对应关系问两组材料中哪些单元表达同一部分内容，近重复问内容重合多少，版本关系还要问哪次编辑继承并改动了哪份材料。三者各有困难。平行文本的句对齐要保序，又要处理增删与一对多，例如译文把一句拆成两句；只靠句长对齐，遇到长度相近的错误译文就无能为力，跨语言句向量让对齐能直接比较语义；近重复检测面对海量文档，逐对比较成本为平方级，MinHash和SimHash用定长签名和分桶把比较降下来。新旧通知高度重合，却只改了一个日期，因此近重复结果必须保留实质差异。

\subsubsection{Gale--Church式句长对齐}
\label{sec:nlp-classification-galechurch}

句长对齐利用两个条件：平行文本大体保序，互译片段的长度成比例。对源片段长度$l_s$和目标片段长度$l_t$，标准化长度差为
\[
  \delta=\frac{l_t-cl_s}{\sqrt{l_s\sigma^2}},
\]
$c$为平均长度比，$\sigma^2$刻画波动。按正态尾概率取匹配代价$-\log[2(1-\Phi(|\delta|))]$，$\Phi$为标准正态分布函数，再加上对齐类型$(1,1)$、$(1,2)$、$(2,1)$、$(2,2)$、$(1,0)$、$(0,1)$的先验罚值\scite{nlp-gale1991alignment}。令$D(i,j)$为源端前$i$句与目标端前$j$句的最小代价，$D(0,0)=0$，
\[
 D(i,j)=\min_{(a,b)}\bigl\{D(i-a,j-b)+C(s_{i-a+1:i},t_{j-b+1:j})\bigr\}.
\]

源文两句长度为$(10,20)$，译文三句长度为$(10,10,10)$，第二句被译成两句。教学设$c=1$、$\sigma^2=10$，$(1,1)$不加罚，$(1,2),(2,1),(2,2)$加1，单边未对应加3；这些罚值人为给定，不是原论文的拟合值。累计表为
\[
 \begin{array}{c|rrrr}
 D(i,j)&0&1&2&3\\ \hline
 0&0&3&6&9\\
 1&3&0&2.1479&5.1479\\
 2&6&2.3935&0.7350&1
 \end{array}
\]
其中$(10,20)$的$\delta=1$，匹配项1.1479，故$D(1,2)=1+1.1479$；$(20,10)$的$\delta=-1/\sqrt2$，匹配项约0.7350；$D(2,1)$取两句合并的$(30,10)$，代价约2.3935。从$D(2,3)=1$回溯，第二源句对目标第2、3句，再由$D(1,1)=0$得到第一句一对一。长度相近不证明翻译正确，摘要式翻译和大段重排都会破坏保序与比例假设。

\subsubsection{跨语言表示匹配与全局对齐}
\label{sec:nlp-classification-crosslingual-align}

长度只是语义对应的弱代理。若有把不同语言映射到同一空间的句编码器，就能直接比较两句的意义。LaBSE是这样的表示型双编码器：在多语言掩码语言模型和翻译语言模型上预训练的BERT，对源句与译句共享参数，取归一化的分类位置向量，用翻译句对训练\scite{nlp-matching-feng2022labse}。训练目标是式~\eqref{eq:nlp-matching-infonce}的双向版本，并在正例分数中减去一个间隔$m$，即加性间隔softmax：要求译句不仅得分最高，还要比批内其他句子至少高出$m$。

用句向量在大语料中挖掘平行句时，直接取余弦最高的目标句会遇到\term{中心点问题}（hubness）：某些泛泛的句子与许多源句都相似，被反复选中。\term{边际打分}（margin-based scoring）把原始余弦除以双方各自近邻的平均相似度\cite{nlp-matching-artetxe2019margin}：
\[
  \operatorname{margin}(x,y)=\frac{\cos(x,y)}
  {\sum_{z\in\mathrm{NN}_k(x)}\frac{\cos(x,z)}{2k}+\sum_{z\in\mathrm{NN}_k(y)}\frac{\cos(y,z)}{2k}},
\]
$\mathrm{NN}_k(x)$是$x$在目标语言中的$k$个最近邻，$\mathrm{NN}_k(y)$是$y$在源语言中的$k$个最近邻。一个句子若与谁都相似，分母就大，分数被压低。设两个源句$x_1,x_2$与三个目标句的余弦为
\[
  \begin{pmatrix}0.80&0.78&0.20\\0.82&0.30&0.75\end{pmatrix},
\]
$y_1$是泛泛的“活动通知”，按余弦两个源句都选它。取$k=2$，$x_1,x_2$的近邻均值为0.79和0.785，$y_1,y_2,y_3$在源端的近邻均值为0.81、0.54、0.475。于是$\operatorname{margin}(x_1,y_1)=0.80/0.80=1.000$，$\operatorname{margin}(x_1,y_2)=0.78/0.665\approx1.173$，$\operatorname{margin}(x_2,y_1)\approx1.028$，$\operatorname{margin}(x_2,y_3)=0.75/0.63\approx1.190$。边际打分让$x_1$选$y_2$、$x_2$选$y_3$，中心句不再独占。这组余弦是教学构造。

挖掘只决定哪些句对可能互译；对已知是平行文本的两篇文档，还要输出保序的全局对齐。Vecalign把Gale--Church的动态规划保留下来，把长度代价换成句向量代价\scite{nlp-matching-thompson2019vecalign}。对候选片段$x,y$（可以是相邻几句的拼接，拼接后重新编码），代价为
\[
  c(x,y)=\frac{(1-\cos(x,y))\,n(x)\,n(y)}
  {\sum_{r=1}^{R}\bigl(1-\cos(x,y_r)\bigr)+\sum_{r=1}^{R}\bigl(1-\cos(x_r,y)\bigr)},
\]
$n(\cdot)$是片段含的句数，$x_r,y_r$是随机抽取的句子。分母用随机句估计“这两个片段与无关句平均相差多少”，消除了不同句子余弦基线的差异；乘以句数是为了抑制无节制的合并。单边未对应使用固定罚值。原文还用逐层下采样的近似动态规划把复杂度降到线性，从而能处理长文档。与逐句独立选最佳不同，全局对齐受保序约束，可以拒绝低质量的配对、让句子保持未对应。评价时应分开候选覆盖率与最终对应正确率：正确译句没进候选是召回问题，进了候选却被约束挤掉是打分或约束问题。

\subsubsection{MinHash：片段集合的近重复检索}
\label{sec:nlp-classification-minhash}

第~\ref{sec:nlp-llm-data-minhash}小节已为预训练语料去重推导了MinHash的碰撞概率等于Jaccard相似度（式~\eqref{eq:nlp-llm-minhash}），以及分带局部敏感哈希的候选概率（式~\eqref{eq:nlp-llm-lsh}）\scite{nlp-broder1997resemblance}。文本组织中用的是同一工具，区别在于目的：语料去重删掉副本即可，这里要找出转载、模板与新旧版本，并保留它们之间的关系。

先决定片段集合怎样构造。连续字符$n$元组适合不分词的中文短通知，词$n$元组对长文更稳；规范化只能去掉空白、编码和全半角等无关差异，不能为提高命中而删掉日期或编号。取$A=\{a,b,c\}$、$B=\{a,b,d\}$，其中$c,d$含不同日期，$J(A,B)=0.5$。短文档片段少，一个日期变化就能让$J$下降很多；长模板共享大量片段，又可能掩盖关键事实变化。当规范要求“新通知主要包含旧通知”时，应计算有方向的包含率$|A\cap B|/|A|$，它与Jaccard回答不同问题。签名只用于召回候选，命中后应计算原集合的相似度并核读正文；近重复只支持合并建议，不能据此删除新版本。

\subsubsection{SimHash与版本差异核验}
\label{sec:nlp-classification-simhash}

MinHash比较集合的交并比例，SimHash近似加权特征向量之间的夹角\scite{nlp-charikar2002similarity}。理想随机超平面方法对随机向量$\bm r_j$记录$\operatorname{sign}(\bm r_j^\top\bm x)$；两个向量夹角为$\theta$时，某一位不同的概率为$\theta/\pi$。工程实现把每个特征哈希成若干正负位，按特征权重逐位累加后取符号，得到一个定长指纹，用汉明距离筛选候选。

给三个特征四位正负码：“读书会”$(+,+,-,-)$、“报名”$(+,-,+,-)$、“周五”$(-,+,+,-)$，权重为2、1、1。逐位和为$(2,2,0,-4)$，约定零取正，指纹为\code{1110}。把“周五”换成码为$(-,-,+,+)$的“周六”，逐位和变为$(2,0,0,-2)$，指纹仍是\code{1110}，汉明距离为0。高权重的共同特征主导了指纹，实质的日期变化被吞没。因此候选命中以后要恢复原文差异，逐项核对实体、数字、否定与条件。新旧版本即使高度重合，也须保存“4月17日”变为“4月18日14:00”的编辑；没有发布时间、活动编号或编辑链记录，仅凭文本无法确定谁先谁后。候选覆盖、重复判定与版本变更核验应分别评价，才能区分漏掉了相近材料还是漏掉了关键修改。

\section{文本集合的组织}
\label{sec:nlp-classification-organization}

前两节的类别和关系都有外部定义。面对一批还没有目录的通知，问题变成怎样发现可供浏览的结构。主题模型给出主题词与文档的混合权重，聚类给出文档簇，关键词给出单篇或一组材料的导航入口。三类结果可以配合使用，但含义不同，评价方式也不同：一个簇名概括了多数成员，不意味着每个成员的每句话都符合它；主题比例、聚类成员分数和关键词得分只在各自模型的定义内有意义。

\subsection{主题发现}
\label{sec:nlp-classification-topics}

\term{主题发现}（topic discovery）从一批文本中寻找反复出现的内容模式，并用词语和代表文档解释。若事先规定类别并有标签，应回到第~\ref{sec:nlp-classification-attributes}节的监督分类。令$x_{d,v}$为文档$d$中词$v$的计数，$K$为主题数，$V$为固定词表。分词、停用词和低频过滤先决定了模型能看见什么：把“无需”删掉以后，词袋方法已无从利用这个否定词。下文固定三个教学词“讲座、报名、场地”，两篇训练文档的计数矩阵为
\begin{equation}
 \label{eq:nlp-classification-topic-counts}
 \bm X=\begin{pmatrix}2&1&0\\0&1&2\end{pmatrix},
\end{equation}
$d_1$围绕讲座与报名，$d_2$围绕报名与场地；另留一篇三词各出现一次的混合文档$d_3=(1,1,1)$，检查新文档怎样使用。

主题发现的困难有五个：结果没有标准答案，难以解释和评价；主题数$K$需要选择；短文本的词袋极其稀疏；新文档的推断有成本；词袋丢掉了词序和角色。方法谱系大致沿这些困难推进。LSA用截断奇异值分解得到潜在维度，但坐标可负、难以解释；NMF加非负约束，成分可读作主题；pLSA给出概率生成解释，用EM估计；LDA加Dirichlet先验，使新文档有后验并缓和短文本过拟合；神经主题模型用推断网络一次前向得到新文档的潜在表示；CTM把上下文句向量接入推断网络，部分弥补词袋丢失的信息；BERTopic改为先聚类句向量、再提取主题词。表~\ref{tab:nlp-classification-topic-methods}对照这些方法。

\begin{table*}[htbp]
 \centering
 \begin{tabularx}{\linewidth}{@{}>{\raggedright\arraybackslash}p{17mm}
   >{\raggedright\arraybackslash}p{30mm}>{\raggedright\arraybackslash}X
   >{\raggedright\arraybackslash}p{30mm}>{\raggedright\arraybackslash}X@{}}
  \toprule
  方法 & 解决的问题 & 核心机制 & 新文档处理 & 新代价与权重含义 \\
  \midrule
  LSA & 同义词分散在不同维度 & 截断SVD，保留前$r$个奇异方向 & 投影到固定右奇异向量 & 坐标可负、符号可翻转，不是概率 \\
  NMF & LSA成分难以解释 & 非负低秩分解，乘性更新 & 固定成分求非负系数 & 尺度不可识别，归一化后也不是后验 \\
  pLSA & 缺少概率解释 & 文档主题比例与主题词分布的混合，EM & 固定词分布局部拟合 & 参数随文档数增长，无文档先验 \\
  LDA & 新文档无定义、短文本过拟合 & Dirichlet先验，折叠Gibbs或变分推断 & 对新文档作后验推断 & 推断迭代代价；主题比例的后验估计 \\
  NVDM、ProdLDA & 逐篇迭代推断太慢 & 推断网络与重参数化，最大化ELBO & 一次前向 & 训练不稳定、主题可能塌缩；高斯坐标或logits组合权重 \\
  CTM & 词袋丢失语序与语境 & 句向量与词袋共同输入推断网络 & 同一编码器与词表 & 依赖句编码器与截断；权重含义同ProdLDA \\
  BERTopic & 短文本词袋稀疏 & 句向量降维、密度聚类、c-TF-IDF选词 & 固定变换与聚类预测接口 & 每篇只归一簇，有离群点；成员分数不是词生成比例 \\
  \bottomrule
 \end{tabularx}
 \caption{主题发现方法的对照。各行基于相同教学材料比较机制，不构成性能排名。潜在坐标、非负系数、主题比例和聚类成员分数即使都落在$[0,1]$，也不能互相替换。}
 \label{tab:nlp-classification-topic-methods}
\end{table*}

\FloatBarrier
\subsubsection{LSA：截断奇异值分解与潜在语义}
\label{sec:nlp-classification-lsa}

“报名”和“登记”在词袋里是两个无关维度，可它们常在相似的文档中出现。\term{潜在语义分析}（latent semantic analysis，LSA）用文档—词矩阵的低秩结构把这些分散的词关联起来：取截断奇异值分解$\bm X\approx\bm U_r\bm\Sigma_r\bm V_r^\top$，$\bm V_r$的每列是一条词空间方向，文档坐标为$\bm X\bm V_r=\bm U_r\bm\Sigma_r$。截断SVD在Frobenius范数下给出最优的秩$r$逼近，非负矩阵分解则在非负约束下求低秩逼近；这两种分解的一般性质可参考\BookChapterRef{chap:dimensionality-reduction}{卷四的降维章}，本节只看它们怎样用于文本。实际使用通常先作TF-IDF加权。

为单独看清分解，本例对式~\eqref{eq:nlp-classification-topic-counts}不加权。两个非零奇异值为$\sqrt6$和2，右方向可取$\bm v_1=(1,1,1)^\top/\sqrt3$、$\bm v_2=(1,0,-1)^\top/\sqrt2$。只保留第一维时，两篇坐标都是$\sqrt3$，重构矩阵两行都是$(1,1,1)$，只剩“都涉及这组词”的共同方向；保留第二维后坐标为$(\sqrt3,\sqrt2)$与$(\sqrt3,-\sqrt2)$，讲座与场地的差别才恢复。新文档沿固定词表加权后右乘$\bm V_r$，不重新拟合；$d_3$投影为$(\sqrt3,0)$。第二方向含负载荷，整列翻转符号也不改变分解，所以它不是“两种主题的概率差”，解释时应分别列出正、负载荷高的词并核读文档。

\subsubsection{NMF：非负成分与主题权重}
\label{sec:nlp-classification-nmf}

LSA的负坐标难以读作“文档含有多少某主题”。若要求每篇文本由非负成分相加解释，可取$\bm X\approx\bm W\bm H$，$\bm W\in\mathbb R_{\geq0}^{D\times K}$为文档系数，$\bm H\in\mathbb R_{\geq0}^{K\times V}$为成分词权重，最小化$\frac12\|\bm X-\bm W\bm H\|_F^2$\scite{lee1999nmf,lee2001nmf}。乘性更新从正初始化出发，交替执行
\[
  \bm W\leftarrow\bm W\odot\frac{\bm X\bm H^\top}{\bm W\bm H\bm H^\top},\qquad
  \bm H\leftarrow\bm H\odot\frac{\bm W^\top\bm X}{\bm W^\top\bm W\bm H},
\]
除法逐元素进行。每一步把元素乘以一个非负比值，非负性自动保持，目标值不增。

对同一$\bm X$取$\bm W=\bigl(\begin{smallmatrix}1&0.5\\0.5&1\end{smallmatrix}\bigr)$、$\bm H=\bigl(\begin{smallmatrix}1&1&0\\0&1&1\end{smallmatrix}\bigr)$，重构两行为$(1,1.5,0.5)$、$(0.5,1.5,1)$，目标值为1.5。固定$\bm H$更新$\bm W$，得$\bm W'=\bigl(\begin{smallmatrix}1.2&0.25\\0.25&1.2\end{smallmatrix}\bigr)$，目标值降至0.905；这只是半轮。两成分的高权重词分别为“讲座、报名”和“报名、场地”，可以读作两个主题。新文档固定$\bm H$求非负最小二乘系数，$d_3$的最优系数为$(2/3,2/3)$。$\bm W$某列乘常数、$\bm H$对应行除以同一常数，重构不变，因此比较系数前必须约定归一化；归一化后的$(1/2,1/2)$是人为消除尺度的结果，不是后验概率。不同初始化可能得到重复成分，非负也不自动带来稀疏。

\subsubsection{pLSA：概率潜在语义分析}
\label{sec:nlp-classification-plsa}

NMF的系数有了可读的符号，却仍没有概率解释，也就无法谈似然与留出预测。\term{概率潜在语义分析}（probabilistic latent semantic analysis，pLSA）把成分约束成分布：文档$d$有主题比例$\theta_{d,k}=p(z=k\mid d)$，主题$k$有词分布$\phi_{k,v}=p(w=v\mid z=k)$，每个词先按文档比例选主题、再按主题选词\cite{pgm-static-directed-hofmann1999}：
\[
 p(v\mid d)=\sum_{k=1}^{K}\theta_{d,k}\phi_{k,v},\qquad
 \ell=\sum_{d,v}x_{d,v}\log\Bigl(\sum_k\theta_{d,k}\phi_{k,v}\Bigr).
\]
求和在对数内部，无法闭式最大化，于是用EM：E步按当前参数计算每个已观测词来自各主题的责任度，M步把责任度当作加权计数重新归一化。EM的一般推导与单调性见第~\ref{sec:nlp-classical-interpolation-em}小节，也可参考\BookChapterRef{chap:pgm-parameter-learning}{卷六的参数学习章}；下面在计数矩阵上完整走一轮。

取两个主题，初始化$\bm\Theta^{(0)}=\bigl(\begin{smallmatrix}0.8&0.2\\0.2&0.8\end{smallmatrix}\bigr)$、$\bm\Phi^{(0)}=\bigl(\begin{smallmatrix}0.6&0.3&0.1\\0.1&0.3&0.6\end{smallmatrix}\bigr)$。E步的责任度为
\[
 r_{d,v,k}=\frac{\theta_{d,k}\phi_{k,v}}{\sum_j\theta_{d,j}\phi_{j,v}}.
\]
$d_1$中“讲座”来自主题1的未归一化质量为$0.8\times0.6=0.48$，来自主题2为$0.02$，责任度为$(0.96,0.04)$；它出现两次，贡献加权计数$(1.92,0.08)$。全部非零项见表~\ref{tab:nlp-classification-plsa-em}。

\begin{table}[htbp]
 \centering
 \begin{tabular}{@{}ccrcccc@{}}
  \toprule
  文档 & 词 & 次数 & $r_1$ & $r_2$ & $x r_1$ & $x r_2$\\
  \midrule
  $d_1$ & 讲座 & 2 & 0.96 & 0.04 & 1.92 & 0.08\\
  $d_1$ & 报名 & 1 & 0.80 & 0.20 & 0.80 & 0.20\\
  $d_2$ & 报名 & 1 & 0.20 & 0.80 & 0.20 & 0.80\\
  $d_2$ & 场地 & 2 & 0.04 & 0.96 & 0.08 & 1.92\\
  \bottomrule
 \end{tabular}
 \caption{pLSA的一次E步。零计数项不贡献加权计数。每行责任度之和为1，加权计数之和等于该词的出现次数。}
 \label{tab:nlp-classification-plsa-em}
\end{table}

M步按文档和主题分别归一化加权计数：
\[
 \theta_{d,k}^{(1)}=\frac{\sum_vx_{d,v}r_{d,v,k}}{\sum_vx_{d,v}},\qquad
 \phi_{k,v}^{(1)}=\frac{\sum_dx_{d,v}r_{d,v,k}}{\sum_{d,u}x_{d,u}r_{d,u,k}}.
\]
$d_1$的主题计数为$(2.72,0.28)$，主题1的词计数为$(1.92,1,0.08)$，总数都是3，于是
$\bm\Theta^{(1)}=\bigl(\begin{smallmatrix}68/75&7/75\\7/75&68/75\end{smallmatrix}\bigr)$，
$\bm\Phi^{(1)}=\bigl(\begin{smallmatrix}16/25&1/3&2/75\\2/75&1/3&16/25\end{smallmatrix}\bigr)$。
对数似然从$-5.1805$升到$-4.3572$。重复两步直到相对增量足够小，并比较多个初始化。新文档没有训练时的$\theta_d$，只能固定$\bm\Phi$、局部迭代拟合它的比例：$d_3$从$(0.5,0.5)$出发，三个词的主题1责任度为$(0.96,0.5,0.04)$，平均仍为0.5。pLSA为每篇训练文档保存自由参数，参数随文档数线性增长，也没有“新文档比例应长什么样”的先验，短文档的比例容易被一两个词支配。

\subsubsection{LDA：潜在狄利克雷分配}
\label{sec:nlp-classification-lda}

\term{潜在狄利克雷分配}（latent Dirichlet allocation，LDA）给文档比例加先验$\bm\theta_d\sim\operatorname{Dir}(\alpha\bm1_K)$，常用的贝叶斯形式还给主题词分布加$\bm\phi_k\sim\operatorname{Dir}(\beta\bm1_V)$\scite{blei2003}。这样新文档的比例有了后验，稀少观测与跨文档共享的先验共同决定结果。生成过程、折叠Gibbs采样的条件分布式~\eqref{eq:nlp-classical-lda-gibbs}及变分推断已在第~\ref{sec:nlp-classical-lda}、\ref{sec:nlp-classical-lda-inference}小节给出，概率图结构可参考\BookChapterRef{chap:pgm-static-directed}{卷六的静态有向模型章}。这里在同一计数矩阵上走通一次采样更新。

把$d_1$的两次“讲座”初始指派给主题1、一次“报名”给主题2，$d_2$的三个词都给主题2；主题词计数为$(2,0,0)$与$(0,2,2)$。取$\alpha=\beta=0.5$，重采样$d_1$的“报名”。先把它从所有计数中移出：文档计数变为$(2,0)$，主题2的词计数变为$(0,1,2)$。按式~\eqref{eq:nlp-classical-lda-gibbs}，
\[
 u_1=(2+0.5)\frac{0+0.5}{2+1.5}=\frac5{14},\qquad
 u_2=(0+0.5)\frac{1+0.5}{3+1.5}=\frac16,
\]
归一化为$(15/22,7/22)$。若均匀随机数为0.6，落入主题1，就把该词加回主题1：文档计数变为$(3,0)$，主题词计数为$(2,1,0)$与$(0,1,2)$。由这一状态估计$\widehat{\bm\theta}_{d_1}=(3.5,0.5)/4=(0.875,0.125)$，$\widehat{\bm\phi}_1=(5/9,1/3,1/9)$。第一个因子偏好文档已经常用的主题，第二个因子偏好能生成当前词的主题；若不先移出旧指派，当前词会给自己增加证据。

这只是一次更新。实际要遍历全部词、多轮采样、丢弃预热阶段，对保留状态取平均；主题编号可以置换，跨运行比较前须先对齐主题。较小的$\alpha$偏好更稀疏的文档比例，较小的$\beta$偏好更集中的主题词分布。先验能缓和短文本的稀疏，却恢复不了词袋丢掉的词序与条件。

\subsubsection{神经主题模型：NVDM与ProdLDA}
\label{sec:nlp-classification-neural-topic}

LDA的新文档推断需要逐篇迭代，主题模型的每次改动也要重新推导采样或变分更新。\term{神经主题模型}（neural topic model）借用变分自编码器的框架：训练一个推断网络，从词袋一次前向就输出潜在分布的参数，再用解码器重构词袋，以证据下界（ELBO）为目标端到端训练。这种做法称为\term{摊销推断}（amortized inference），因为推断的代价被分摊到网络参数里。变分自编码器与重参数化的推导可参考\BookChapterRef{chap:pgm-generative}{卷六的概率生成模型章}。

NVDM是这一路线的直接形式\scite{nlp-miao2016nvdm}。推断网络把计数$\bm x\in\mathbb R^{V}$映射为$\bm\mu(\bm x),\log\bm\sigma^2(\bm x)\in\mathbb R^{K}$，取$\bm h=\bm\mu+\bm\sigma\odot\bm\epsilon$、$\bm\epsilon\sim\mathcal N(\bm0,\bm I)$，解码器以$\operatorname{softmax}(\bm R\bm h+\bm b)$给出词分布$\bm p(\bm h)$，$\bm R\in\mathbb R^{V\times K}$。训练最大化
\[
 \mathcal L=\mathbb E_q\Bigl[\sum_vx_v\log p_v(\bm h)\Bigr]
 -D_{\mathrm{KL}}\bigl(q(\bm h\mid\bm x)\,\|\,\mathcal N(\bm0,\bm I)\bigr).
\]
对$d_1=(2,1,0)$，设推断网络输出$\bm\mu=(0.5,0)$、$\bm\sigma^2=(1,1)$，取一份教学噪声$\bm\epsilon=\bm0$，解码logits为$(0.5,0,-0.5)$，词概率约$(0.5065,0.3072,0.1863)$。重构负对数似然约2.5408，KL项为$\frac12\times0.5^2=0.125$，负ELBO的单样本估计约2.6658。对logits的重构梯度为$N\bm p-\bm x\approx(-0.48,-0.08,0.56)$，$N=3$为词数，更新会提高“讲座、报名”的概率。NVDM的$\bm h$可正可负、不加和为1，不是主题比例。

ProdLDA让潜在表示回到单纯形上，并改变词分布的组合方式\scite{nlp-srivastava2017prodlda}。Dirichlet分布无法直接重参数化，ProdLDA在softmax基下用Laplace近似把Dirichlet先验换成logistic正态分布：推断网络仍输出高斯参数，采样后取$\bm\theta=\operatorname{softmax}(\bm h)$。LDA的词分布是主题词分布的算术混合$\sum_k\theta_k\bm\phi_k$；ProdLDA先在logits上组合主题，再归一化：
\[
 p_v=\frac{\exp\bigl(\sum_k\theta_kb_{k,v}\bigr)}{\sum_u\exp\bigl(\sum_k\theta_kb_{k,u}\bigr)}.
\]
若取$b_{k,v}=\log\phi_{k,v}$，它就是归一化加权几何平均$\prod_k\phi_{k,v}^{\theta_k}$，这就是\term{专家乘积}（product of experts）。主题词分布取$(0.6,0.3,0.1)$与$(0.1,0.3,0.6)$、权重各0.5时，LDA混合得$(0.35,0.30,0.35)$，ProdLDA得约$(0.3101,0.3798,0.3101)$：乘积压低了只受单个主题偏好的词，两个主题都给中等质量的“报名”反而排第一。原文认为这种组合让主题词更集中，这里的数值只演示机制。

两种模型都把新文档推断变成一次前向，也继承了变分自编码器的问题：KL项可能把后验压向先验，使部分主题从不被使用；不同运行的主题也可能大量重复。验收时应检查主题词重复、各主题的使用比例和跨运行稳定性。

\subsubsection{CTM：上下文文档表示与主题推断}
\label{sec:nlp-classification-ctm}

词袋推断网络无法区分词相同而次序不同的材料，短文本的词袋又太稀疏。CTM（contextualized topic model）把SBERT句向量接入ProdLDA的推断网络。CombinedTM把句向量投影后与词袋拼接作为输入\cite{nlp-intro-bianchi2021}；ZeroShotTM只用句向量作输入，因此可以用多语言句编码器为其他语言的文档推断主题\scite{nlp-bianchi2021ctm}。两种配置的解码器都重构词袋，主题仍由词分布解释。

把“甲通知乙”与“乙通知甲”简化为相同词袋$(1,1,1)$，句向量分别为$(1,0)$与$(0,1)$。若推断均值相应为$(1,0)$与$(0,1)$，零噪声下主题权重为$(0.7311,0.2689)$与$(0.2689,0.7311)$。相同词袋得到了不同的推断输入；但训练目标仍是重构词袋，若这些词的分布不需要区分角色，学到的主题也可能相同。新文档必须使用同一句编码器、投影和词表；长文档被句编码器截断后，遗漏的条件不会因拼接而恢复。

\subsubsection{BERTopic：表示聚类与c-TF-IDF}
\label{sec:nlp-classification-bertopic}

前面的方法都从词袋出发。BERTopic反过来，先用句向量把文档分组，再从每组的词袋中挑代表词\scite{nlp-grootendorst2022bertopic}。流程分三步：SBERT等编码器得到文档向量；UMAP降维后由HDBSCAN寻找密度簇，并把不属于任何簇的文档标为离群；把每个簇的文档合并成一篇，计算类别级TF-IDF（c-TF-IDF）选出代表词。向量决定谁和谁近，词袋只用于解释已经得到的簇，两个环节可以分别调整。

设聚类输出把$d_1$与补充的讲座通知$(2,0,0)$归入簇A，$d_2$归入簇B，$d_3$为离群。这组指派是教学给定的，不是由三维计数运行UMAP得到。簇计数为$A=(4,1,0)$、$B=(0,1,2)$，平均簇长$\overline L_c=4$。c-TF-IDF为
\[
 w_{c,v}=\frac{x_{c,v}}{\sum_ux_{c,u}}\log\Bigl(1+\frac{\overline L_c}{\sum_{c'}x_{c',v}}\Bigr),
\]
分母是词$v$在所有簇中的总次数，不是出现该词的簇数。三词的逆频率项为$(\log2,\log3,\log3)$，得$\bm w_A\approx(0.5545,0.2197,0)$、$\bm w_B\approx(0,0.3662,0.7324)$，两簇可分别命名为“讲座、报名”与“场地、报名”。合并两簇时须重新计算。BERTopic每篇文档只归一个簇，混合文档$d_3$要么成离群、要么被硬分，这与LDA的主题混合不同；簇成员分数反映密度结构，不是词生成比例。

\subsubsection{主题数选择、解释与评价}
\label{sec:nlp-classification-topic-evaluation}

主题词看起来合理，只说明通过了一次人工印象检查。自动指标从三方面补充：\term{主题连贯性}（topic coherence）看高权重词在参考语料中是否常共同出现，主题多样性看不同主题是否反复使用同一组词，跨运行稳定性先对齐主题再比较词分布。连贯性常用归一化点互信息。若100个固定窗口中“讲座”出现20次、“报名”30次、共同出现10次，
\[
 \operatorname{NPMI}(u,v)=\frac{\log[p(u,v)/(p(u)p(v))]}{-\log p(u,v)}
 =\frac{\log(0.1/0.06)}{-\log0.1}\approx0.2218.
\]
它在$[-1,1]$内，1表示总是同现。共现为零时须按评价协议平滑；经常同现的也可能只是模板页脚。人工评价用词语侵入和文档侵入测试：把一个异类词混入主题词表或把一篇异类文档混入代表文档，看读者能否找出\scite{nlp-organization-chang2009tea}。原文的实验显示，留出似然更好的模型，侵入测试的结果不一定更好。

留出似然或每词困惑度只适用于概率生成路线，且须固定词表与推断协议；BERTopic等聚类路线没有同名的似然。选择主题数时，固定语料与预处理，在若干$K$下比较重复主题、连贯性、多样性、稳定性与用途；密度聚类则记录簇数和离群率怎样随参数变化。数量过少会合并讲座与场地，过多会把同一模板拆成多个主题。

\subsection{文档聚类}
\label{sec:nlp-classification-clustering}

文档聚类把表示相近的材料放在一起，“相近”由表示、归一化和距离共同决定。文本聚类的困难在于：TF-IDF向量高维稀疏，句向量维数低却各向异性；簇数事先未知；混合文档和离群文档很常见；簇的解释必须回到真实文档。同主题与同版本还需要不同的身份约束：主题聚类允许不同日期的同类活动放在一起，版本聚类必须先区分活动编号。K-means、层次聚类与密度聚类的定义、目标与收敛性质可参考\BookChapterRef{chap:clustering}{卷四的聚类章}；本小节只讨论文本任务中的选择。

本小节把四份通知的二维表示固定为$A=(1,0)$、$B=(0.8,0.2)$、$C=(0,1)$、$D=(0.2,0.8)$，另以$E=(0.5,0.5)$代表混合材料。向量均为教学构造。

\subsubsection{文本表示下的聚类选择}
\label{sec:nlp-classification-cluster-choice}

先选表示，再选距离。TF-IDF向量的范数随文档长度和词频增长：同一份通知重复两遍得到$2A=(2,0)$，它与$A$的欧氏距离为1，比$A$与$B$的距离0.2828还大，余弦却等于1。文本的“内容相同”通常对应方向相同，因此稀疏词向量一般用余弦。对单位向量，
\[
  \|\bm x-\bm y\|_2^2=2-2\bm x^\top\bm y,
\]
欧氏距离的排序与余弦一致；先把所有文档$L_2$归一化，再使用任何基于欧氏距离的算法，就等于在余弦下聚类。句向量常已归一化，但其分布往往集中在一个狭窄锥体内，绝对的余弦值不宜跨编码器比较，阈值须在当前编码器上验证。

再选是否允许软分配。硬划分给每篇文档一个簇；在连续表示上拟合高斯混合$p(\bm x)=\sum_k\pi_k\mathcal N(\bm x;\bm\mu_k,\bm\Sigma_k)$，则由成分后验给出软分配。取两成分等先验，均值$(0.9,0.1)$与$(0.1,0.9)$，共享协方差$0.08\bm I$。新通知$F=(0.7,0.3)$到两中心的平方距离为0.08与0.72，第一成分后验为$\sigma\bigl((0.72-0.08)/(2\times0.08)\bigr)=\sigma(4)\approx0.9820$。$E$到两中心等距，后验各0.5。但后验各半不等于内容各半：一篇离两簇都很远的文档同样可能后验各半，因此还要检查绝对密度或另设拒绝条件。这个后验说的是“文档向量属于哪个成分”，与LDA中“每个词来自哪个主题”的比例不是同一对象。

最后是评价。\term{轮廓系数}（silhouette coefficient）$s_i=(b_i-a_i)/\max(a_i,b_i)$比较样本到同簇的平均距离$a_i$与到最近其他簇的平均距离$b_i$；有参考分组时可补充\term{调整兰德指数}（adjusted Rand index, ARI）与归一化互信息（normalized mutual information, NMI）：前者统计样本对是否被两种划分同样地分在一起，并扣除随机划分下的期望值；后者把两种划分之间的互信息归一化到$[0,1]$。内部指标与簇名都不能替代对代表文档的核读，也不能替代新文档分配的实际验证。

\subsubsection{球面K-means文档聚类}
\label{sec:nlp-classification-kmeans}

普通K-means以平方欧氏距离分配、以簇内均值更新中心，中心的范数会随成员变化，与余弦口径不符。球面K-means把文档和中心都限制在单位球面上\cite{nlp-organization-dhillon2001concept}：目标是最大化$\sum_k\sum_{\bm x\in\mathcal C_k}\bm x^\top\bm c_k$，约束$\|\bm c_k\|_2=1$。分配步把每篇文档分给余弦最大的中心；更新步取簇内向量之和再归一化，$\bm c_k=\sum_{\bm x\in\mathcal C_k}\bm x/\|\sum_{\bm x\in\mathcal C_k}\bm x\|_2$，它在单位约束下使簇内余弦之和最大。由前一节的恒等式，这等价于对归一化文档做平方欧氏距离K-means，只是中心多了一步归一化。

对$A,B,C,D$取初始中心$(1,0)$与$(0,1)$。归一化后$B$变为$(0.9701,0.2425)$，分到第一簇；$D$对称地分到第二簇。第一簇中心为$A$与归一化$B$之和再归一化，约为$(0.9925,0.1222)$，而普通K-means的均值是$(0.9,0.1)$。重新分配不变，迭代停止。$E$到两个中心的余弦相同，需要事先规定并列规则，或把它标为低确定性，而不是硬称它只属于一个主题。中心是抽象方向，不是任何原文，解释簇时要取离中心最近的真实文档。球面K-means适合高维稀疏的TF-IDF，每轮代价与非零元个数成正比；它仍要预设$K$，结果依赖初始化，空簇需要按预设规则重启。

\subsubsection{密度聚类与层次聚类}
\label{sec:nlp-classification-density}

K-means类方法要求预设簇数，并把每篇文档都分进某个簇。密度聚类与层次聚类从不同方向放松这两点。DBSCAN以半径$\varepsilon$和最少邻域点数$\mathrm{MinPts}$定义核心点，沿核心点的密度可达关系恢复簇，邻域稀疏的点标为噪声\scite{ester1996}。取欧氏距离、$\varepsilon=0.3$、$\mathrm{MinPts}=2$（邻域计入自身），$A,B$相距0.2828，$C,D$亦然，各成一簇；$E$到最近的$B,D$都约0.4243，邻域只有自身，成为噪声。把$\varepsilon$增至0.45，$E$同时连到$B$和$D$，成为核心点，两簇被连成一簇。这个变化只反映密度连通的变化，不说明两组通知突然变成了同一主题。

单一半径难以同时适应密集与稀疏的区域。HDBSCAN用核心距离（到第$\mathrm{MinPts}$个近邻的距离，本例邻域计入自身）定义互达距离$d_{\mathrm{mr}}(a,b)=\max\{\mathrm{core}(a),\mathrm{core}(b),d(a,b)\}$，在其上构造所有半径下的簇层次，再按簇在层次中存在的稳定性选出最终簇\scite{clust-campello2013}。本例$\mathrm{MinPts}=2$时，$A$的核心距离约0.2828，$E$约0.4243，$E$参与的边都处在较稀疏的层级，因此$E$最后才并入、或保持为离群。这正是BERTopic选用它的原因：短文本中大量文档不属于任何清晰主题，强行归簇会污染代表词。

凝聚层次聚类则不预设簇数，而是保留整个合并过程：从每篇一簇开始，反复合并最近的两簇，记录合并距离形成树。簇间距离的定义决定结果。$\{A,B\}$与$\{C,D\}$之间，单链接取最近点对$B,D$的0.8485，全链接取最远点对$A,C$的1.4142，平均链接取四个点对的平均约1.1314；在距离1.0处截断，单链接已合成一簇，另两种仍是两簇。加入桥接文档$E$后，单链接更容易沿局部近邻串成一个大簇，这就是链式效应。层次树表示合并顺序，不自动对应“活动通知—讲座—机器学习”这样的业务层级；经典算法也没有自然的新文档分配方式，重新聚类可能改变已有目录。三类方法的共同验收是：同时报告覆盖率与噪声比例，在非噪声子集上计算内部指标时注明排除了谁，并核读代表文档。

\subsection{关键词与关键短语组织}
\label{sec:nlp-classification-keyphrases}

文档分组后，用户仍需要简短的入口来找到材料。\term{关键词抽取}（keyword extraction）从原文中选出词或跨度，\term{关键短语生成}（keyphrase generation）允许输出原文未出现的概括。单篇关键词服务当前材料，主题代表词解释潜在分布，簇代表词解释文档组，同一个词在不同接口下的得分并不相同。

这一任务有四个困难。候选生成决定上限，漏掉“报名截止时间”这个候选，后面的排序无法补回；短语边界须完整，“机器学习读书会”不能只剩“读书会”；所选短语彼此冗余会浪费导航位置；参考关键短语中常有原文没出现的概括。方法谱系依次回应这些困难：TF-IDF依赖参考语料的区分度；RAKE与YAKE只用单篇局部统计，不需要语料；TextRank在共现图上传播重要性；EmbedRank用语义相似排序并以MMR去冗余；序列标注与跨度方法从标注中学习边界规范；生成式方法可以产生原文未出现的短语。表~\ref{tab:nlp-classification-keyphrase-methods}汇总这些方法。

\begin{table*}[htbp]
 \centering
 \begin{tabularx}{\linewidth}{@{}>{\raggedright\arraybackslash}p{24mm}
   >{\raggedright\arraybackslash}p{32mm}>{\raggedright\arraybackslash}X
   >{\raggedright\arraybackslash}p{40mm}@{}}
  \toprule
  方法 & 解决的问题 & 核心机制 & 新代价 \\
  \midrule
  TF-IDF & 区分常见词与局部重要词 & 单篇词频乘参考语料的逆文档频率 & 需要参考语料；短语须另造候选 \\
  RAKE、YAKE & 无参考语料 & 停用词切分加共现度；位置、频率、分散度等局部特征 & 依赖停用词表与分词；启发式权重 \\
  TextRank & 局部频率忽略词间关系 & 共现图上的PageRank & 窗口与词性过滤敏感 \\
  EmbedRank & 统计量不看语义，结果冗余 & 候选与文档嵌入的余弦，MMR去冗余 & 依赖句编码器；需调$\lambda$ \\
  序列标注、跨度 & 边界规范难以手写 & BIO加CRF，或跨度打分 & 需要标注；无法输出原文外短语 \\
  CopyRNN、catSeq & 原文未出现的概括 & 编码器—解码器加复制；单条或整组生成 & 生成内容须核验；整组生成引入顺序 \\
  \bottomrule
 \end{tabularx}
 \caption{关键短语方法的对照。前五类只能输出原文中的跨度，评价应把原文内与原文外的参考短语分开统计。}
 \label{tab:nlp-classification-keyphrase-methods}
\end{table*}

\FloatBarrier
\subsubsection{TF-IDF：词频与跨文档区分度}
\label{sec:nlp-classification-keyword-tfidf}

TF-IDF用单篇词频与参考语料的文档频率区分常见词和局部重要词。采用$\operatorname{idf}(v)=\log((N+1)/(\operatorname{df}_v+1))+1$：参考集有4篇，“读书会”只出现在1篇、在当前通知出现2次，权重为$2(\log(5/2)+1)\approx3.8326$；“通知”4篇都出现、当前出现3次，权重为3。前者次数少，却因区分度更高排在前面。参考语料的来源必须说明，过度的低频过滤会删掉活动编号和新术语。短语候选可由词性模式或连续$n$元组生成，再按整短语的文档频率计分，或按规定聚合组成词的权重；高分的“机器”和“报名”若不相邻，不能拼成一条抽取短语。

\subsubsection{RAKE与YAKE：单篇文本的局部统计}
\label{sec:nlp-classification-rake}

没有参考语料时，只能从当前文档内部找信号。RAKE先用停用词和分隔符把文本切成候选短语，再用词在候选中的共现计分\scite{nlp-rose2010rake}。词频$f(w)$为出现次数，共现度$\deg(w)$按包含自身的口径计：词每出现在一个长为$r$的候选中，度加$r$。词分数为$\deg(w)/f(w)$，短语分数为其词分数之和。固定候选“机器／学习／读书会”“报名／信息”“机器／学习”，则“机器、学习”频次各2、度各5，分数各2.5；“读书会”分数3；“报名、信息”分数各2。三个候选得分为8、4、5，排序为完整活动短语、“机器学习”、“报名信息”。长短语天然得分高，这是$\deg/f$偏好长候选的结果；中文的分词与停用词表直接决定候选边界。

YAKE改用更多局部特征，并让分数越小越好\scite{nlp-campos2020yake}。对单词$w$，$C$为大写特征，$F=f/(\overline f+s_f)$为相对词频（$\overline f,s_f$为非停用词词频的均值与标准差），$D$为出现句数占比，$P=\log\log(3+\text{所在句序号的中位数})$为位置特征，$R=1+(L+R_{\!c})f/f_{\max}$刻画左右邻词的分散程度，其中$L,R_{\!c}$为左、右邻词的种类数与邻接次数之比。组合为
\[
 H(w)=\frac{PR}{C+F/R+D/R},\qquad
 H(c)=\frac{\prod_{w\in c}H(w)}{f(c)\bigl[1+\sum_{w\in c}H(w)\bigr]}.
\]
出现越早$P$越小；频繁且分布在多句中使分母增大；邻词越杂乱$R$越大，说明它更像虚词。设某词$f=2$、$f_{\max}=4$、$\overline f=2$、$s_f=1$，在4句中的第0、1句出现，$L=R_{\!c}=0.5$，无大小写信号，则$C=0$、$F=2/3$、$D=1/2$、$P=\log\log3.5$、$R=1.5$，$H\approx0.4346$；它与$H=0.3$的另一词组成只出现一次的候选，得$H(c)\approx0.0752$。中文没有大小写，$C$恒为零。两种方法都无需外部语料，代价是标题重复与模板开头可能被过度奖励。

\subsubsection{TextRank：共现图与关键词排序}
\label{sec:nlp-classification-textrank}

局部统计把词孤立地计分。TextRank认为，与许多重要词共现的词本身也重要：候选词为节点，在固定窗口内共现的词之间连边，再做PageRank迭代\scite{nlp-mihalcea2004textrank}。对$M$个节点，取阻尼$d\in(0,1)$，
\[
 r_i^{(t+1)}=\frac{1-d}{M}+d\sum_{j:\,w_{j,i}>0}\frac{w_{j,i}}{\sum_kw_{j,k}}r_j^{(t)}.
\]
取“机器、学习、读书会”三个节点，只有无向边“机器—学习”“学习—读书会”。初值各$1/3$、$d=0.85$，一轮后两端各约0.1917，中间约0.6167，总和仍为1；迭代到相邻向量的$L_1$差小于容差才得到稳定排名。选定高分词后，按原文中的连续出现恢复“机器学习读书会”，不能把不相邻的高分词随意连接。窗口过大会让通用词连到许多对象，词性过滤也可能删去关键动词。第~\ref{chap:nlp-text-generation}章的抽取式摘要把节点换成句子，复用同一图排序。

\subsubsection{EmbedRank：语义表示与MMR选择}
\label{sec:nlp-classification-embedrank}

统计方法不看语义，选出的前几个短语也常彼此重复。EmbedRank用词性规则取名词短语候选，把候选和文档都编码为句向量，按余弦排序；EmbedRank++再用\term{最大边际相关性}（maximal marginal relevance，MMR）平衡相关与重复\scite{nlp-bennani2018embedrank}。已选集合为$S$时，从未选候选中取
\[
 \operatorname{MMR}(c)=\lambda s(c,x)-(1-\lambda)\max_{c'\in S}s(c,c')
\]
的最大者，$x$为文档。为便于手算，这里直接用原始余弦，原文还对相似度作了规范化。候选A“机器学习读书会”、B“读书会”、C“报名信息”与文档的相似度为0.90、0.88、0.75，A与B相似度0.95，A与C为0.40。先选A；$\lambda=0.7$时B得0.331、C得0.405，第二项选C；$\lambda=0.9$时B得0.697、C得0.635，又改选B。$\lambda$表达导航对重复的容忍度，没有普遍正确的取值。“报名信息”虽提示了报名事项，却没有保留截止时刻，还要核查所选候选是否覆盖关键事项。

\subsubsection{序列标注与跨度式关键短语抽取}
\label{sec:nlp-classification-keyphrase-spans}

有参考关键短语时，边界规范可以从标注中学习。把参考短语在原文中的出现位置转成BIO标签，就成为第~\ref{sec:nlp-analysis-paradigm-tagging}小节的序列标注；枚举有限长度跨度并打分，就成为第~\ref{sec:nlp-analysis-paradigm-span}小节的跨度方法。转换前须先规定：同一短语多次出现是否都标，嵌套短语是否都保留，原文中不存在的参考短语从抽取训练中排除。

对“机器／学习／读书会／报名／截止”，若参考为前3词与后2词，BIO序列为$(B,I,I,B,I)$。独立逐词分类若把第4词判为$O$、第5词判为$I$，就出现非法的$O\to I$；CRF的转移约束可以禁止它。跨度模型可给$(1,3)$与$(4,5)$打分0.92与0.86，在验证阈值0.8下恢复两段，并映射回字符区间。若同时需要“机器学习”和“机器学习读书会”，平坦BIO无法表达嵌套，需改用跨度集合。评价应分别报告候选覆盖、完整边界F1与去重后的短语F1，以区分没找到位置、边界不全与选错短语。

\subsubsection{生成式关键短语预测：CopyRNN与One2Seq}
\label{sec:nlp-classification-copyrnn}

抽取方法输出不了原文中没有的概括，而科技文献的作者关键词中有相当一部分不在原文里出现。CopyRNN把关键短语预测改写为序列生成：双向GRU编码文档，带注意力的GRU解码器每次生成一条短语，并用复制机制让解码器能直接输出原文中的低频词和未登录词\scite{nlp-meng2017copyrnn}。训练数据按“一篇文档—一条参考短语”拆成多个样本，因此称为One2One。复制与生成的一般形式见式~\eqref{eq:nlp-analysis-copy}，那里用门控概率$p_{\mathrm{gen}}$混合两个分别归一化的分布；CopyRNN沿用CopyNet的写法，把生成与复制的未归一化质量放进同一个分母：
\[
 p_t(w)=\frac{\mathbf1[w\in V]\exp g_t(w)+\sum_{j:x_j=w}\exp c_t(j)}
 {\sum_{v\in V}\exp g_t(v)+\sum_{j=1}^{|x|}\exp c_t(j)},
\]
$g_t(v)$为生成logit，$c_t(j)$为复制原文第$j$个位置的logit。两种写法的输出空间都是词表与原文词的并集，同一个词在原文多处出现时复制质量相加。

设某步词表只展示“报名、信息、结束符”，生成质量为$(2,1,1)$；原文三个位置为“海棠楼、报名、报名”，复制质量为$(1,1,2)$。共同分母为8，$p(\text{报名})=(2+1+2)/8=0.625$，不在词表中的“海棠楼”也可经复制得到$1/8$。训练目标是参考短语及结束符的负对数似然：三步条件概率为$(0.625,0.5,0.8)$时，序列概率为0.25，损失约1.3863。推断时用束搜索得到多条完成的序列（束搜索的一般做法见第~\ref{sec:nlp-generation-beam}小节），去重后取前若干条。One2One的问题是输出数量需要人为截断，例如固定取前5条，而且各条短语独立生成，彼此冗余。

One2Seq把一篇文档的全部参考短语用分隔符拼成一个目标序列，例如“机器学习读书会$\langle\mathrm{sep}\rangle$报名截止$\langle\mathrm{sep}\rangle$海棠楼$\langle\mathrm{eos}\rangle$”，由同一个带复制的编码器—解码器生成；catSeq是这种训练方式下的模型\scite{nlp-organization-yuan2020catseq}。这样做有两个直接收益：生成后一条短语时能看到已生成的短语，可以避免重复；模型自己决定何时输出结束符，从而决定短语数量。原文发现，贪心解码一条序列就能输出数量可变的短语，并相应提出按模型实际输出计算的F1@M。代价是把无序的集合变成了有序的序列：同一组短语的不同排列给出不同的训练目标，常用的约定是原文中出现的短语按首次出现位置排在前，未出现的短语排在后；顺序选择会影响结果。

大语言模型也可以直接按提示生成关键短语，提示中给出数量、格式与少样本示例即可（第~\ref{sec:nlp-inference-icl-shots}小节）。无论哪种生成方式，原文内的短语都可以用字符串匹配回查，原文外的概括则需单独核验是否得到原文支持。评价时应把原文内与原文外的参考短语分开统计，并声明精确匹配还是词形规范化：能复制地点不证明它最有用，能生成流畅的概括也不证明原文支持这个概括。


\section*{文献说明}

本章的原始资料按支持的内容分为三类：模型结构、标签与数据规范、评价口径。

分类模型方面，fastText与TextCNN分别给出平均嵌入与卷积分类的原始设计\cite{nlp-joulin2017fasttext,nlp-kim2014textcnn}，层次注意力网络给出词级与句级两层注意力的结构\cite{nlp-yang2016han}；XML-CNN与AttentionXML给出极端多标签的瓶颈层、动态池化、标签注意力与概率标签树\cite{nlp-liu2017xmlcnn,nlp-you2019attentionxml}。方面级情感分析依次参考ATAE-LSTM、IAN、BERT-pair、BARTABSA、Span-ASTE与GTS\cite{nlp-wang2016atae,nlp-ma2017ian,nlp-sun2019bertpair,nlp-yan2021bartabsa,nlp-xu2021spanaste,nlp-wu2020gts}；未见目标立场检测参考TGA-Net及其VAST数据，以及TOAD\cite{nlp-allaway2020vast,nlp-allaway2021toad}。来源判定参考DetectGPT、Fast-DetectGPT与绿名单水印\cite{nlp-mitchell2023detectgpt,nlp-bao2024fastdetectgpt,nlp-kirchenbauer2023watermark}，生成参与的四类划分依据Creator-Editor研究\cite{nlp-li2026creator}。

匹配模型方面，表示型与交互型架构参考DSSM、ARC与MatchPyramid\cite{rs-h10,nlp-matching-hu2014arc,nlp-matching-pang2016matchpyramid}，ESIM给出上下文编码与局部推断组合\cite{nlp-chen2017esim}；句向量与对比学习参考Sentence-BERT与SimCSE\cite{nlp-reimers2019sbert,nlp-gao2021simcse}；检索式匹配参考DPR、ColBERT与Poly-encoder\cite{nlp-karpukhin2020dpr,nlp-khattab2020colbert,nlp-matching-humeau2020polyencoder}，难负例挖掘与交叉编码器蒸馏参考ANCE、RocketQA与Margin-MSE\cite{rs-s03,nlp-matching-qu2021rocketqa,nlp-matching-hofstatter2020margin}；跨语言对齐参考LaBSE、边际打分与Vecalign\cite{nlp-matching-feng2022labse,nlp-matching-artetxe2019margin,nlp-matching-thompson2019vecalign}；近重复检测参考MinHash与SimHash的原始分析\cite{nlp-broder1997resemblance,nlp-charikar2002similarity}。组织方法参考NVDM、ProdLDA、CTM与BERTopic\cite{nlp-miao2016nvdm,nlp-srivastava2017prodlda,nlp-intro-bianchi2021,nlp-bianchi2021ctm,nlp-grootendorst2022bertopic}，球面K-means参考Dhillon与Modha的概念分解\cite{nlp-organization-dhillon2001concept}，主题的人工可解释性评价参考词语与主题侵入实验\cite{nlp-organization-chang2009tea}；关键词参考YAKE、EmbedRank、CopyRNN与catSeq\cite{nlp-campos2020yake,nlp-bennani2018embedrank,nlp-meng2017copyrnn,nlp-organization-yuan2020catseq}。

数据与评价规范方面，情感任务的粒度划分参考SemEval-2017 Task 4\cite{nlp-rosenthal2017sentiment}，立场的目标类型与未见目标设定参考Bi-STANCE\cite{nlp-zhao2025bistance}；语义相似度、推断与复述难例分别参考STS、SNLI与MultiNLI、PAWS\cite{nlp-intro-cer2017,nlp-bowman2015snli,nlp-intro-williams2018,nlp-intro-zhang2019paws}。各节的数值例子均为教学构造；引用原文结果时只保留说明机制所需的个别数字，不能据此对方法作统一的性能排名。

\section{本章小结}
\label{sec:nlp-classification-summary}

回到章首的挑战。分类从词袋线性模型走到fastText与TextCNN，再到层次注意力网络、预训练编码器与提示式分类，每一步回答前一步留下的问题：同义词不共享统计、局部短语需要枚举、长文档的关键句被稀释、上下文交互要从少量标注中学出、新标签需要重新训练。多标签任务用逐标签sigmoid与验证集阈值，标签依赖可由分类器链建模；标签数达到数十万时，瓶颈层、标签注意力与概率标签树把输出层的计算降下来，评价改用P@$k$与nDCG@$k$。情感、立场与来源判定各需一个额外条件：方面级情感分析从注意力模型走到方面、观点与极性的统一抽取；立场以目标为条件，未见目标要靠相近目标或与目标无关的表示迁移；来源判定从监督分类、困惑度、DetectGPT到Fast-DetectGPT逐步降低零样本检测的成本，水印则以生成方配合换来有统计保证的检测。这些统计信号都不构成作者身份的证明。

匹配按两段文本何时交互在精度与可预计算之间取舍：表示型适合大规模召回，交互型适合小候选集上的细判，后期交互介于二者之间。相似、等价、蕴含与相关是不同的关系，对称与否、候选集合怎样构成都要随任务说明，一种分数不能替代另一种。对齐要保留未配对与一对多，近重复检测还要核对日期、实体和版本来源。

没有预设类别时，主题模型、聚类与关键词给出三种组织结果：主题比例、非负系数、潜在坐标和簇成员分数各有模型内的含义，不应统称为主题概率；关键词把材料变成可浏览的入口，候选、排序与去重都需核查。到这里，本部分处理的输出都是标签、配对、原文片段或组织结构。下一部分转向生成新文本，输出不再受原文位置约束，需要另外检查它保留了什么、改写了什么，以及是否加入了原文没有依据的内容。
